\documentclass[11pt]{article}

% ============================================================
% Packages
% ============================================================
\usepackage[margin=1in]{geometry}
\usepackage{amsmath,amssymb,mathtools,amsthm}
\usepackage{enumitem}
\usepackage{hyperref}
\usepackage{xcolor}
\usepackage{microtype}

\usepackage{listings}
\usepackage{xcolor}

\lstdefinestyle{mypython}{
    language=Python,
    basicstyle=\ttfamily\small,
    keywordstyle=\color{blue},
    commentstyle=\color{teal},
    stringstyle=\color{red},
    showstringspaces=false,
    breaklines=true,
    breakatwhitespace=false,
    columns=fullflexible,
    keepspaces=true,
    frame=single,
    tabsize=4
}

% ============================================================
% Macros
% ============================================================
\newcommand{\R}{\mathbb{R}}
\newcommand{\Qbb}{\mathbb{Q}}
\newcommand{\E}{\mathbb{E}}
\newcommand{\dd}{\,\mathrm{d}}

% ============================================================
% Environments
% ============================================================
\theoremstyle{plain}
\newtheorem{theorem}{Theorem}[section]
\newtheorem{proposition}[theorem]{Proposition}

\theoremstyle{definition}
\newtheorem{definition}[theorem]{Definition}
\newtheorem{example}[theorem]{Example}

\theoremstyle{remark}
\newtheorem{remark}[theorem]{Remark}

% ============================================================
% Title
% ============================================================
\title{Option Pricing --- Lecture Notes}
\author{Nathan Sauldubois}
\date{MSFE6233 --- Fixed-Income Derivatives Pricing}

\begin{document}
% ============================================================
\section{European option: pricing, hedging, and model risk}
% ============================================================
\subsection{The contract and the general objective}

We begin with the simplest and most classical derivative contract: a
European call option.

\paragraph{Contract.}
A European call option with maturity \(T\) and strike \(K\) gives its
holder the right, but not the obligation, to buy the underlying asset at
time \(T\) for the price \(K\).

Its payoff is
\[
(S_T-K)^+.
\]

This is the natural benchmark contract for several reasons:
\begin{itemize}[leftmargin=1.4em]
    \item the payoff depends only on the terminal value \(S_T\),
    \item the contract is liquid in many markets,
    \item its price can be computed explicitly in the Black--Scholes
    model,
    \item and its hedge can be derived in a fully explicit way.
\end{itemize}

\paragraph{General objective of this section.}
Our goal is not only to compute the price of the call, but also to
understand how a pricing model generates a hedging strategy.

More precisely, we want to explain the following chain of ideas:
\begin{quote}
a model specifies the dynamics of the underlying asset; from these
dynamics, one derives a pricing formula and a pricing PDE; from the PDE,
one extracts the delta; and this delta defines the dynamic hedging
strategy.
\end{quote}

In the first step, we work in the Black--Scholes model, which provides a
simple and fully explicit benchmark.

Later in the section, we will revisit the same contract under a local
volatility model and compare the corresponding hedging strategies.

\subsection{A benchmark model: Black--Scholes}

We now recall the Black--Scholes framework, which will serve as the
benchmark model throughout this first example.

\subsubsection{Risk-neutral dynamics}

We assume that the market contains:
\begin{itemize}[leftmargin=1.4em]
    \item a risky asset \(S\),
    \item and a money market account \(B\).
\end{itemize}

The money market account satisfies
\[
\mathrm{d}B_t=rB_t\,\mathrm{d}t,
\qquad
B_0=1,
\]
so that
\[
B_t=e^{rt},
\]
where \(r\) is the constant risk-free interest rate.

\paragraph{Black--Scholes assumption under the risk-neutral measure.}
Under the risk-neutral measure \(\Qbb\), the asset price is assumed to
follow the stochastic differential equation
\[
\mathrm{d}S_t=rS_t\,\mathrm{d}t+\sigma S_t\,\mathrm{d}W_t,
\]
where:
\begin{itemize}[leftmargin=1.4em]
    \item \(r\) is the constant interest rate,
    \item \(\sigma>0\) is the constant volatility,
    \item \(W\) is a Brownian motion under \(\Qbb\).
\end{itemize}

\paragraph{Interpretation.}
The risk-neutral dynamics say that, under \(\Qbb\), the expected rate of
return of the asset is the risk-free rate \(r\), not the physical
expected return of the stock.

This is the key simplification of arbitrage-free pricing:
once we move to the risk-neutral measure, derivative prices are obtained
by discounting future payoffs at the risk-free rate.

\paragraph{Explicit solution.}
The Black--Scholes SDE admits the explicit solution
\[
S_T
=
S_t\exp\!\Big(
\big(r-\tfrac12\sigma^2\big)(T-t)+\sigma(W_T-W_t)
\Big).
\]

Thus the stock price is lognormal under \(\Qbb\), and the log-return
\[
\log\Big(\frac{S_T}{S_t}\Big)
\]
is Gaussian with mean
\[
\big(r-\tfrac12\sigma^2\big)(T-t)
\]
and variance
\[
\sigma^2(T-t).
\]

This explicit structure is what makes the Black--Scholes model so
tractable.

\subsubsection{Risk-neutral pricing formula}

Let \(u(t,s)\) denote the price at time \(t\) of the European call when
the current stock price is \(s\).

\paragraph{Risk-neutral valuation.}
Under the usual assumptions of no arbitrage and market completeness, the
time-\(t\) price of the call is given by
\[
u(t,s)
=
\E^\Qbb\!\left[
e^{-r(T-t)}(S_T-K)^+
\;\middle|\;
S_t=s
\right].
\]

Equivalently, along the stochastic path,
\[
u(t,S_t)
=
\E^\Qbb\!\left[
e^{-r(T-t)}(S_T-K)^+
\;\middle|\;
\mathcal F_t
\right].
\]

\paragraph{Interpretation.}
This formula says that the option price is the discounted conditional
expectation of its payoff under the risk-neutral measure.

The remarkable point is that pricing does not involve the investor's
risk preferences explicitly.
Everything is encoded in:
\begin{itemize}[leftmargin=1.4em]
    \item the risk-neutral dynamics of the stock,
    \item and the discounting at the risk-free rate.
\end{itemize}

\paragraph{Why this formula is natural.}
If the discounted stock price
\[
e^{-rt}S_t
\]
is a martingale under \(\Qbb\), then by the same no-arbitrage logic the
discounted price of any self-financing replicable claim must also be a
martingale.
This leads directly to the risk-neutral expectation formula.

\paragraph{Dependence on the current state.}
Because the Black--Scholes model is Markovian, the price depends only on
the current time \(t\) and the current stock price \(s\), so it can be
written as a deterministic function
\[
u=u(t,s).
\]

This Markovian structure is what allows us to derive a PDE.

\subsubsection{Pricing PDE}

We now derive the PDE satisfied by the price function \(u(t,s)\).

\paragraph{It\^o's formula.}
Assume that \(u\) is sufficiently smooth.
Applying It\^o's formula to the process \(u(t,S_t)\), we obtain
\[
\mathrm{d}u(t,S_t)
=
\partial_t u(t,S_t)\,\mathrm{d}t
+
\partial_s u(t,S_t)\,\mathrm{d}S_t
+
\frac12 \partial_{ss}u(t,S_t)\,\mathrm{d}\langle S\rangle_t.
\]

Since
\[
\mathrm{d}S_t=rS_t\,\mathrm{d}t+\sigma S_t\,\mathrm{d}W_t
\]
and
\[
\mathrm{d}\langle S\rangle_t=\sigma^2 S_t^2\,\mathrm{d}t,
\]
this becomes
\[
\mathrm{d}u(t,S_t)
=
\Big(
\partial_t u
+
rS_t\,\partial_s u
+
\frac12 \sigma^2 S_t^2\,\partial_{ss}u
\Big)\mathrm{d}t
+
\sigma S_t\,\partial_s u\,\mathrm{d}W_t.
\]

\paragraph{Discounted martingale argument.}
Under the risk-neutral measure, the discounted price process
\[
e^{-rt}u(t,S_t)
\]
must be a martingale.
Applying It\^o's product rule, its drift must vanish.
This yields the Black--Scholes PDE:
\[
\boxed{
\partial_t u(t,s)
+
rs\,\partial_s u(t,s)
+
\frac12 \sigma^2 s^2\,\partial_{ss}u(t,s)
-
r\,u(t,s)
=
0.
}
\]

The terminal condition is given by the payoff:
\[
\boxed{
u(T,s)=(s-K)^+.
}
\]

\paragraph{Interpretation of the PDE.}
This equation summarizes the whole pricing problem:
\begin{itemize}[leftmargin=1.4em]
    \item the terminal condition encodes the contract,
    \item the differential operator encodes the model dynamics,
    \item and the discount term \(-ru\) reflects the time value of money.
\end{itemize}

\paragraph{Alternative derivation by replication.}
The same PDE can also be obtained by constructing a hedged portfolio and
imposing that it earns the risk-free rate.
We will do this explicitly in the next subsection, because it leads
directly to the delta hedging strategy.

\subsubsection{Closed-form price of the European call}

One of the main strengths of the Black--Scholes model is that the PDE
can be solved explicitly for the European call.

The resulting price is
\[
\boxed{
u(t,s)=s\,N(d_1)-Ke^{-r(T-t)}N(d_2),
}
\]
where
\[
d_1
=
\frac{\log(s/K)+\big(r+\tfrac12\sigma^2\big)(T-t)}
{\sigma\sqrt{T-t}},
\qquad
d_2=d_1-\sigma\sqrt{T-t}.
\]

Here \(N(\cdot)\) denotes the cumulative distribution function of a
standard normal random variable:
\[
N(x)=\mathbb{P}(Z\le x),
\qquad
Z\sim\mathcal N(0,1).
\]

\paragraph{Interpretation of the formula.}
The Black--Scholes price has a very natural structure:
\begin{itemize}[leftmargin=1.4em]
    \item the term \(sN(d_1)\) represents the weighted contribution of
    the stock,
    \item the term \(Ke^{-r(T-t)}N(d_2)\) represents the weighted present
    value of the strike.
\end{itemize}

\paragraph{Dependence on the parameters.}
The formula makes it possible to read immediately how the call price
depends on:
\begin{itemize}[leftmargin=1.4em]
    \item the current stock price \(s\),
    \item the strike \(K\),
    \item the interest rate \(r\),
    \item the volatility \(\sigma\),
    \item and the time to maturity \(T-t\).
\end{itemize}

\paragraph{Why this formula matters for us.}
In this section, the closed-form price is useful for two reasons:
\begin{itemize}[leftmargin=1.4em]
    \item it provides an explicit benchmark price,
    \item and it allows us to compute the hedge ratio explicitly.
\end{itemize}

This explicit delta will be the reference point for our later
comparison with local volatility hedging.

\subsection{Delta hedging in the Black--Scholes model}

We now turn to the hedging problem.
The key question is the following:

\begin{quote}
If we sell the option, how should we dynamically trade in the stock and
the bank account in order to replicate its payoff?
\end{quote}

This leads to the notion of delta hedging.

\subsubsection{Self-financing portfolio}

Suppose that we want to replicate the option by trading dynamically in:
\begin{itemize}[leftmargin=1.4em]
    \item the stock \(S_t\),
    \item the money market account \(B_t\).
\end{itemize}

Let:
\begin{itemize}[leftmargin=1.4em]
    \item \(\Delta_t\) be the number of shares held in the stock,
    \item \(\beta_t\) be the amount invested in the bank account.
\end{itemize}

The value of the hedging portfolio is then
\[
\Pi_t=\Delta_t S_t+\beta_t B_t.
\]

\paragraph{Self-financing condition.}
We assume that the strategy is self-financing, meaning that changes in
portfolio value come only from gains and losses on the assets already
held, not from external cash injections or withdrawals.

Therefore,
\[
\mathrm{d}\Pi_t=\Delta_t\,\mathrm{d}S_t+\beta_t\,\mathrm{d}B_t.
\]

Since
\[
\mathrm{d}B_t=rB_t\,\mathrm{d}t,
\]
we have
\[
\mathrm{d}\Pi_t=\Delta_t\,\mathrm{d}S_t+r\beta_t B_t\,\mathrm{d}t.
\]

Using
\[
\Pi_t=\Delta_t S_t+\beta_t B_t,
\]
this can also be written as
\[
\mathrm{d}\Pi_t
=
\Delta_t\,\mathrm{d}S_t+r(\Pi_t-\Delta_t S_t)\,\mathrm{d}t.
\]

\paragraph{Replication objective.}
We want the portfolio to replicate the option, that is,
\[
\Pi_t=u(t,S_t)
\qquad\text{for all }t.
\]
In particular, at maturity, we want
\[
\Pi_T=(S_T-K)^+.
\]

Thus the problem is to choose \(\Delta_t\) and \(\beta_t\) so that the
portfolio value always coincides with the option price.

\subsubsection{Derivation of the hedge ratio}

We now compare the dynamics of the option price and the dynamics of the
self-financing portfolio.

\paragraph{Dynamics of the option price.}
From It\^o's formula, we already know that
\[
\mathrm{d}u(t,S_t)
=
\Big(
\partial_t u
+
rS_t\,\partial_s u
+
\frac12 \sigma^2 S_t^2\,\partial_{ss}u
\Big)\mathrm{d}t
+
\sigma S_t\,\partial_s u\,\mathrm{d}W_t.
\]

\paragraph{Dynamics of the portfolio.}
On the other hand,
\[
\mathrm{d}\Pi_t
=
\Delta_t\,\mathrm{d}S_t+r(\Pi_t-\Delta_t S_t)\,\mathrm{d}t.
\]
Since
\[
\mathrm{d}S_t=rS_t\,\mathrm{d}t+\sigma S_t\,\mathrm{d}W_t,
\]
we obtain
\[
\mathrm{d}\Pi_t
=
r\Pi_t\,\mathrm{d}t+\Delta_t \sigma S_t\,\mathrm{d}W_t.
\]

\paragraph{Matching the diffusion terms.}
If the portfolio is to replicate the option, then we must have
\[
\Pi_t=u(t,S_t),
\]
and therefore the Brownian terms in the two dynamics must coincide.

This gives
\[
\Delta_t \sigma S_t
=
\sigma S_t\,\partial_s u(t,S_t).
\]

Assuming \(\sigma S_t\neq 0\), we deduce
\[
\boxed{
\Delta_t=\partial_s u(t,S_t).
}
\]

This is the hedge ratio.

\paragraph{Interpretation.}
The delta hedge ratio is the sensitivity of the option price with
respect to the underlying asset price.
It tells us how many units of the stock must be held at time \(t\) in
order to offset the first-order exposure of the option to moves in the
underlying.

\paragraph{Matching the drift terms.}
Once the diffusion terms are matched, the drift terms also coincide
provided that \(u\) solves the Black--Scholes PDE.
Thus the PDE and the replicating strategy are two sides of the same
argument:
\begin{itemize}[leftmargin=1.4em]
    \item the hedge ratio removes the random part,
    \item the PDE ensures consistency of the deterministic part.
\end{itemize}

\subsubsection{The Black--Scholes delta}

Since we have an explicit formula for the call price, we can compute the
delta explicitly.

Recall that
\[
u(t,s)=sN(d_1)-Ke^{-r(T-t)}N(d_2).
\]

Differentiating with respect to \(s\), one obtains the classical
Black--Scholes delta:
\[
\boxed{
\Delta^{BS}(t,s)=\partial_s u(t,s)=N(d_1).
}
\]

Therefore, along the path of the stock,
\[
\boxed{
\Delta_t^{BS}=N(d_1(t,S_t)).
}
\]

\paragraph{Meaning of the formula.}
This quantity lies between \(0\) and \(1\), which is consistent with the
fact that the call option becomes more and more stock-like as it moves
deeper into the money.

More precisely:
\begin{itemize}[leftmargin=1.4em]
    \item if the option is deep out of the money, \(\Delta^{BS}\) is
    close to \(0\),
    \item if the option is deep in the money, \(\Delta^{BS}\) is close
    to \(1\),
    \item near the money, \(\Delta^{BS}\) takes intermediate values.
\end{itemize}

\paragraph{Dynamic nature of the hedge.}
The delta is not constant.
It changes with time and with the level of the stock price.
Therefore, the replicating strategy must be rebalanced dynamically.

This is why one speaks of \emph{dynamic hedging}.

\paragraph{Discrete rebalancing in practice.}
In the ideal Black--Scholes theory, rebalancing is continuous in time.
In practice, however, hedging is done on a discrete grid of dates.
This generates hedging error even if the model is correct.
Later in the section, we will study this issue numerically.

\subsubsection{Interpretation of the replication argument}

The Black--Scholes replication argument is one of the most important
ideas in mathematical finance.

\paragraph{Main message.}
The option price is not introduced as an arbitrary discounted expected
value.
Rather, it is the unique value of a self-financing trading strategy that
replicates the payoff.

Thus pricing and hedging are intrinsically linked:
\begin{itemize}[leftmargin=1.4em]
    \item the pricing formula tells us the value of the contract,
    \item the hedge ratio tells us how to trade dynamically in order to
    replicate it.
\end{itemize}

\paragraph{Why this is powerful.}
The argument shows that, in a complete model, the value of the option is
fully determined by no arbitrage.
If two different prices were proposed for the same replicable payoff,
one could construct an arbitrage.

\paragraph{The role of the model.}
The whole construction depends on the assumed dynamics of the
underlying.
Once the model is fixed, the PDE and the hedge ratio follow from it.

This is the key point for the rest of the course:
\begin{quote}
different models generate different prices, different PDEs, and
different hedging strategies.
\end{quote}

\paragraph{Why this benchmark matters.}
The Black--Scholes model gives us:
\begin{itemize}[leftmargin=1.4em]
    \item a closed-form price,
    \item an explicit PDE,
    \item an explicit delta hedge,
    \item and a fully transparent replication argument.
\end{itemize}

This makes it the perfect benchmark against which more sophisticated
models can be compared.

\paragraph{Transition to the next subsection.}
In the next step, we keep the same contract, namely the European call,
but we change the model.
Instead of assuming a constant volatility, we use a local volatility
surface calibrated from liquid market option prices.

This will allow us to compare:
\begin{itemize}[leftmargin=1.4em]
    \item a simple benchmark hedge, namely the Black--Scholes delta,
    \item and a market-consistent hedge, namely the local-volatility
    delta.
\end{itemize}

\subsection{A market-consistent extension: the local volatility model}

We now keep the same contract, namely the European call, but replace the
Black--Scholes assumption of constant volatility by a more flexible
model calibrated from market vanilla prices.

The main idea is the following:
\begin{quote}
even when the price of a liquid European option is already observed in
the market, one still needs a dynamic model for the underlying in order
to compute Greeks and construct a hedging strategy.
\end{quote}

This is one of the main motivations for the local volatility model.

\subsubsection{Why go beyond constant volatility?}

The Black--Scholes model is a natural benchmark, but it relies on the
assumption that the volatility coefficient \(\sigma\) is constant.
This implies that all European option prices at a given date should be
consistent with a single volatility parameter.

In practice, this is not what one observes.
Market option prices, when translated into Black--Scholes implied
volatilities, exhibit clear patterns:
\begin{itemize}[leftmargin=1.4em]
    \item implied volatility depends on the strike,
    \item implied volatility depends on the maturity,
    \item and the resulting implied-volatility surface is not flat.
\end{itemize}

In equity markets, this often appears as a downward skew:
low strikes tend to have higher implied volatilities than high strikes.
Thus the assumption of constant volatility is too rigid to match the
observed cross-section of liquid European option prices.

\paragraph{A first consequence for hedging.}
If the true market smile is not flat, then the Black--Scholes delta is
not necessarily the most market-consistent hedge ratio.
Indeed, the option price reacts not only to the stock level, but also to
the way volatility varies with strike and maturity.

Thus, even for a European call, the choice of pricing model affects the
computed delta and therefore the hedging strategy.

\paragraph{What we want from a richer model.}
We would like a model that:
\begin{itemize}[leftmargin=1.4em]
    \item remains relatively tractable,
    \item is still Markovian in the stock price,
    \item reproduces the full current surface of liquid European option
    prices,
    \item and provides a dynamic framework for hedging.
\end{itemize}

This leads naturally to the local volatility approach.

\subsubsection{Calibration from liquid European option prices}

The key practical feature of the local volatility model is that it is
not introduced in isolation.
Rather, it is built from the observed market prices of liquid European
options.

\paragraph{Observed market data.}
At a given date, the market provides prices
\[
C^{\mathrm{mkt}}(T,K)
\]
for a family of liquid European call options indexed by maturity \(T\)
and strike \(K\).

Equivalently, one may think of the market as providing the implied
volatility surface
\[
\sigma_{\mathrm{imp}}(T,K),
\]
since Black--Scholes gives a one-to-one relation between call prices and
implied volatilities.

\paragraph{What is the objective?}
The purpose of the local volatility model is to construct a deterministic
function
\[
\sigma_{\mathrm{loc}}(t,s)
\]
such that, when the stock follows the corresponding diffusion, the model
reproduces the observed market prices of liquid European options.

Thus the calibration problem is:
\[
\text{find }\sigma_{\mathrm{loc}}(t,s)\text{ such that }
C^{\mathrm{model}}(T,K)=C^{\mathrm{mkt}}(T,K)
\]
for the relevant maturities and strikes.

\paragraph{Why this is useful even for a liquid option.}
One might ask:
if the call price is already quoted in the market, why do we need a
model at all?

The answer is that the model is needed not primarily to reprice the
liquid option itself, but to provide:
\begin{itemize}[leftmargin=1.4em]
    \item a coherent risk-neutral dynamics for the underlying,
    \item a pricing PDE,
    \item and, most importantly for us here, a delta hedge
    \[
    \Delta_t=\partial_s u(t,S_t).
    \]
\end{itemize}

Thus, in the local volatility setting, market prices are used as
\emph{inputs} to construct a dynamic hedging model.

\paragraph{Static fit versus dynamic use.}
This is a very important point.
Local volatility is calibrated from today's market vanilla prices, so it
is excellent at fitting the current implied-volatility surface.
Once calibrated, it generates a full dynamic model for the stock price,
and this model can then be used:
\begin{itemize}[leftmargin=1.4em]
    \item to compute Greeks,
    \item to hedge,
    \item and to price less liquid products.
\end{itemize}

\subsubsection{Dupire's local volatility formula}

We now recall the central formula that links liquid call prices to the
local volatility surface.

Suppose that the market call-price surface
\[
C(T,K)
\]
is sufficiently smooth in maturity and strike.
Then Dupire's formula formally gives the local variance as
\[
\boxed{
\sigma_{\mathrm{loc}}(T,K)^2
=
\frac{
\partial_T C(T,K)+rK\,\partial_K C(T,K)
}{
\frac12 K^2\,\partial_{KK}C(T,K)
}.
}
\]

Equivalently, the call price surface satisfies the forward Dupire PDE
\[
\partial_T C(T,K)
=
\frac12 \sigma_{\mathrm{loc}}(T,K)^2K^2\,\partial_{KK}C(T,K)
-rK\,\partial_K C(T,K).
\]

\paragraph{Interpretation.}
This formula is remarkable because it expresses the local volatility
surface directly in terms of observable market quantities:
\begin{itemize}[leftmargin=1.4em]
    \item the derivative of call prices with respect to maturity,
    \item and the curvature of call prices with respect to strike.
\end{itemize}

Thus the market surface of vanilla options is turned into a local
diffusion coefficient for the stock.

\paragraph{What does the formula really mean?}
The quantity
\[
\sigma_{\mathrm{loc}}(T,K)
\]
should not be confused with the implied volatility.
The implied volatility is the Black--Scholes volatility that matches one
single option price.
The local volatility is instead the instantaneous diffusion coefficient
of the stock process in the model calibrated to the whole surface.

\paragraph{Practical point.}
In applications, one does not observe a perfectly smooth surface
\(C(T,K)\), but only a finite grid of strikes and maturities.
Hence interpolation and numerical differentiation are needed in order to
apply Dupire's formula.
This is one of the numerical difficulties of the local volatility
approach.

\paragraph{What matters for us here.}
For the purposes of this section, the key point is the following:
once the market call surface is given, Dupire's formula provides a way
to infer a state-dependent volatility function
\[
\sigma_{\mathrm{loc}}(t,s),
\]
and this function then determines the model used to compute the hedge.

\subsubsection{Risk-neutral dynamics under local volatility}

Once the local volatility surface has been constructed, the stock price
is modeled under the risk-neutral measure by the diffusion
\[
\boxed{
\mathrm{d}S_t
=
rS_t\,\mathrm{d}t
+
\sigma_{\mathrm{loc}}(t,S_t)S_t\,\mathrm{d}W_t.
}
\]

Thus the drift remains the risk-free rate, exactly as in the
Black--Scholes model, but the diffusion coefficient is no longer
constant.

\paragraph{Comparison with Black--Scholes.}
Under Black--Scholes,
\[
\mathrm{d}S_t=rS_t\,\mathrm{d}t+\sigma S_t\,\mathrm{d}W_t.
\]
Under local volatility, this becomes
\[
\mathrm{d}S_t=rS_t\,\mathrm{d}t+\sigma_{\mathrm{loc}}(t,S_t)S_t\,\mathrm{d}W_t.
\]

So the model remains:
\begin{itemize}[leftmargin=1.4em]
    \item one-dimensional,
    \item Markovian,
    \item and risk-neutral.
\end{itemize}

But it is now flexible enough to reflect the current option surface.

\paragraph{Log-price dynamics.}
Applying It\^o's formula to \(\log S_t\), we get
\[
\mathrm{d}\log S_t
=
\Big(
r-\frac12 \sigma_{\mathrm{loc}}(t,S_t)^2
\Big)\mathrm{d}t
+
\sigma_{\mathrm{loc}}(t,S_t)\,\mathrm{d}W_t.
\]

Unlike Black--Scholes, the log-price is no longer Gaussian in general,
since the volatility depends on the current state.

\paragraph{Discounted dynamics.}
If we define the discounted stock
\[
\widetilde S_t=e^{-rt}S_t,
\]
then
\[
\mathrm{d}\widetilde S_t
=
e^{-rt}\sigma_{\mathrm{loc}}(t,S_t)S_t\,\mathrm{d}W_t,
\]
so the discounted stock remains a local martingale under the
risk-neutral measure.

\paragraph{Main practical interpretation.}
The local volatility model is therefore a market-consistent extension of
Black--Scholes:
\begin{itemize}[leftmargin=1.4em]
    \item the drift is unchanged,
    \item the dynamics remain relatively simple,
    \item but the volatility coefficient is chosen so as to reflect the
    current market smile.
\end{itemize}

\subsection{Pricing and hedging in the local volatility model}

We now use the local volatility dynamics to define the option price and
the corresponding delta hedge.

\subsubsection{Risk-neutral pricing formula}

Let \(u(t,s)\) denote the price of the same European call option under
the local volatility model.

\paragraph{Risk-neutral valuation.}
Since the model is still written under the risk-neutral measure, the
pricing formula remains
\[
\boxed{
u(t,s)
=
\E^\Qbb\!\left[
e^{-r(T-t)}(S_T-K)^+
\;\middle|\;
S_t=s
\right],
}
\]
where now the process \(S\) follows the local volatility dynamics
\[
\mathrm{d}S_t
=
rS_t\,\mathrm{d}t
+
\sigma_{\mathrm{loc}}(t,S_t)S_t\,\mathrm{d}W_t.
\]

\paragraph{Same contract, different model.}
This is an important point:
the payoff has not changed,
\[
(S_T-K)^+,
\]
but the law of \(S_T\) under the pricing measure is no longer the
Black--Scholes lognormal law.
The price function \(u(t,s)\) is therefore different from the
Black--Scholes price function, even though the contract is the same.

\paragraph{Markovian structure.}
Since the local volatility model remains one-factor and Markovian, the
price still depends only on \((t,s)\).
This allows us to derive a PDE very similar to Black--Scholes, but with
state-dependent volatility.

\subsubsection{Pricing PDE under local volatility}

Assume that \(u\) is sufficiently smooth.
Applying It\^o's formula to \(u(t,S_t)\), we obtain
\[
\mathrm{d}u(t,S_t)
=
\partial_t u\,\mathrm{d}t
+
\partial_s u\,\mathrm{d}S_t
+
\frac12 \partial_{ss}u\,\mathrm{d}\langle S\rangle_t.
\]

Since
\[
\mathrm{d}S_t
=
rS_t\,\mathrm{d}t
+
\sigma_{\mathrm{loc}}(t,S_t)S_t\,\mathrm{d}W_t,
\]
and
\[
\mathrm{d}\langle S\rangle_t
=
\sigma_{\mathrm{loc}}(t,S_t)^2S_t^2\,\mathrm{d}t,
\]
this becomes
\[
\mathrm{d}u(t,S_t)
=
\Big(
\partial_t u
+
rS_t\partial_s u
+
\frac12 \sigma_{\mathrm{loc}}(t,S_t)^2S_t^2\partial_{ss}u
\Big)\mathrm{d}t
+
\sigma_{\mathrm{loc}}(t,S_t)S_t\partial_s u\,\mathrm{d}W_t.
\]

As before, under the risk-neutral measure the discounted option price
must be a martingale.
Therefore the drift of
\[
e^{-rt}u(t,S_t)
\]
must vanish, which yields the PDE
\[
\boxed{
\partial_t u(t,s)
+
rs\,\partial_s u(t,s)
+
\frac12 \sigma_{\mathrm{loc}}(t,s)^2 s^2\,\partial_{ss}u(t,s)
-r\,u(t,s)
=
0,
}
\]
with terminal condition
\[
\boxed{
u(T,s)=(s-K)^+.
}
\]

\paragraph{Interpretation.}
This PDE has exactly the same structure as the Black--Scholes PDE, but
with the constant coefficient \(\sigma\) replaced by the local
volatility function \(\sigma_{\mathrm{loc}}(t,s)\).

Thus the local volatility model preserves the PDE logic of
Black--Scholes while incorporating the market smile through the diffusion
coefficient.

\paragraph{No closed-form formula in general.}
Unlike in the Black--Scholes model, one usually does not have a closed
formula for \(u(t,s)\) in the local volatility setting.
Therefore the price and the delta must generally be computed
numerically, for instance by finite-difference methods or Monte Carlo.

\subsubsection{Delta hedge under local volatility}

The hedging logic is the same as in the Black--Scholes model:
we construct a self-financing portfolio made of stock and cash, and we
choose the stock position so as to match the diffusion term of the
option price.

\paragraph{Self-financing portfolio.}
Let
\[
\Pi_t=\Delta_t S_t+\beta_t B_t
\]
be the value of a self-financing hedging portfolio, where \(B_t=e^{rt}\)
is the money market account.

Then
\[
\mathrm{d}\Pi_t=\Delta_t\,\mathrm{d}S_t+\beta_t\,\mathrm{d}B_t.
\]

As in the Black--Scholes case, this can be rewritten as
\[
\mathrm{d}\Pi_t
=
r\Pi_t\,\mathrm{d}t
+
\Delta_t \sigma_{\mathrm{loc}}(t,S_t)S_t\,\mathrm{d}W_t.
\]

\paragraph{Match the Brownian terms.}
If we want
\[
\Pi_t=u(t,S_t),
\]
then we compare this with the dynamics of \(u(t,S_t)\):
\[
\mathrm{d}u(t,S_t)
=
\cdots
+
\sigma_{\mathrm{loc}}(t,S_t)S_t\,\partial_s u(t,S_t)\,\mathrm{d}W_t.
\]

Matching the diffusion terms yields
\[
\Delta_t \sigma_{\mathrm{loc}}(t,S_t)S_t
=
\sigma_{\mathrm{loc}}(t,S_t)S_t\,\partial_s u(t,S_t),
\]
and therefore
\[
\boxed{
\Delta_t^{LV}=\partial_s u(t,S_t).
}
\]

\paragraph{Same formal formula, different quantity.}
Formally, the delta hedge has exactly the same definition as before:
it is still the derivative of the option price with respect to the spot.

However, the price function \(u\) is no longer the Black--Scholes one.
Hence:
\[
\Delta_t^{LV}\neq \Delta_t^{BS}
\]
in general.

This is one of the most important lessons of this section:
\begin{quote}
the hedge ratio is model-dependent because the price function is
model-dependent.
\end{quote}

\paragraph{Practical implication.}
In the local volatility model, the delta is usually computed numerically
from the PDE solution.
Thus, compared with Black--Scholes:
\begin{itemize}[leftmargin=1.4em]
    \item the conceptual hedging argument is unchanged,
    \item but the hedge ratio is no longer explicit,
    \item and its numerical value reflects the market smile encoded in
    \(\sigma_{\mathrm{loc}}(t,s)\).
\end{itemize}

\subsubsection{Comparison with the Black--Scholes hedge}

We now compare the two delta hedges.

\paragraph{Black--Scholes hedge.}
Under Black--Scholes, the hedge ratio is
\[
\Delta_t^{BS}=N(d_1(t,S_t)).
\]
It is explicit and depends only on:
\begin{itemize}[leftmargin=1.4em]
    \item the stock level,
    \item the strike,
    \item the maturity,
    \item the constant volatility,
    \item and the interest rate.
\end{itemize}

\paragraph{Local volatility hedge.}
Under local volatility, the hedge ratio is
\[
\Delta_t^{LV}=\partial_s u(t,S_t),
\]
where \(u\) solves the local volatility PDE.
This delta depends on the entire local volatility surface
\[
\sigma_{\mathrm{loc}}(t,s),
\]
and therefore on the observed market smile.

\paragraph{Conceptual difference.}
The Black--Scholes delta is a hedge produced by a model that ignores the
smile.
The local-volatility delta is a hedge produced by a model calibrated to
the current smile.

Thus the local-volatility hedge is more \emph{market-consistent} in the
sense that it uses information extracted from liquid market option
prices.

\paragraph{Important practical nuance.}
For a liquid European call, the market already provides the option
price.
So the role of the local volatility model is not mainly to compute a new
price for this liquid contract.
Its role is rather to provide:
\begin{itemize}[leftmargin=1.4em]
    \item a dynamic model for the stock,
    \item a coherent PDE,
    \item and therefore a hedge ratio consistent with the implied
    volatility surface.
\end{itemize}

\paragraph{What we expect qualitatively.}
If the market smile is pronounced, then the local volatility delta may
differ significantly from the Black--Scholes delta, especially:
\begin{itemize}[leftmargin=1.4em]
    \item for strikes far from the money,
    \item for longer maturities,
    \item or when the local volatility surface varies strongly with the
    spot.
\end{itemize}

\paragraph{Transition to the numerical section.}
At this stage, we have two candidate hedging strategies for the same
European call:
\begin{itemize}[leftmargin=1.4em]
    \item the Black--Scholes delta hedge,
    \item the local-volatility delta hedge.
\end{itemize}

In the next part, we explain how to compute these quantities
numerically, rebalance the corresponding self-financing portfolios in
discrete time, and compare their hedging performance when the true
underlying dynamics are stochastic volatility rather than either of the
two hedging models.


\subsection{Numerical implementation}

We now explain how to implement the pricing and hedging strategies
numerically.

The overall workflow is the following:
\begin{itemize}[leftmargin=1.4em]
    \item solve the pricing PDE backward in time,
    \item recover the option delta from the numerical solution,
    \item simulate stock paths,
    \item rebalance the hedging portfolio on a discrete time grid,
    \item compare the resulting hedging P\&L across models.
\end{itemize}

In the Black--Scholes case, the price and the delta are explicit, so the
numerical burden is light.
In the local volatility case, the price and the delta must typically be
computed on a grid.

\subsubsection{Finite-difference scheme for the pricing PDE}

We start from the local volatility pricing PDE
\[
\partial_t u(t,s)
+
rs\,\partial_s u(t,s)
+
\frac12 \sigma_{\mathrm{loc}}(t,s)^2 s^2\,\partial_{ss}u(t,s)
-r\,u(t,s)
=
0,
\]
with terminal condition
\[
u(T,s)=(s-K)^+.
\]

\paragraph{Space-time grid.}
Fix:
\begin{itemize}[leftmargin=1.4em]
    \item a terminal time \(T\),
    \item a maximum stock value \(S_{\max}\),
    \item integers \(N_t,N_s\ge 1\).
\end{itemize}

We define the time and space steps
\[
\Delta t=\frac{T}{N_t},
\qquad
\Delta s=\frac{S_{\max}}{N_s},
\]
and the grid
\[
t_n=n\Delta t,\qquad n=0,\dots,N_t,
\]
\[
s_i=i\Delta s,\qquad i=0,\dots,N_s.
\]

We denote by
\[
u_i^n \approx u(t_n,s_i)
\]
the numerical approximation of the option price at the grid point
\((t_n,s_i)\).

\paragraph{Terminal and boundary conditions.}
At maturity,
\[
u_i^{N_t}=(s_i-K)^+.
\]

For the spatial boundaries, a simple and standard choice for a call is:
\[
u_0^n=0,
\qquad
u_{N_s}^n=S_{\max}-Ke^{-r(T-t_n)}.
\]

The first condition reflects the fact that a call is worthless when the
stock is zero.
The second one is the usual linear asymptotic behavior of the call for
large stock prices.

\paragraph{Backward explicit scheme.}
A simple explicit backward scheme is obtained by writing
\[
\partial_t u(t_n,s_i)
\approx
\frac{u_i^{n+1}-u_i^n}{\Delta t},
\]
\[
\partial_s u(t_n,s_i)
\approx
\frac{u_{i+1}^{n+1}-u_{i-1}^{n+1}}{2\Delta s},
\]
\[
\partial_{ss}u(t_n,s_i)
\approx
\frac{u_{i+1}^{n+1}-2u_i^{n+1}+u_{i-1}^{n+1}}{\Delta s^2}.
\]

Substituting these approximations into the PDE and solving for \(u_i^n\)
gives
\[
u_i^n
=
u_i^{n+1}
+
\Delta t
\Bigg[
rs_i\frac{u_{i+1}^{n+1}-u_{i-1}^{n+1}}{2\Delta s}
+
\frac12 \sigma_{\mathrm{loc}}(t_n,s_i)^2 s_i^2
\frac{u_{i+1}^{n+1}-2u_i^{n+1}+u_{i-1}^{n+1}}{\Delta s^2}
-r u_i^{n+1}
\Bigg].
\]

This allows us to compute the solution backward from \(n=N_t-1\) down
to \(n=0\).

\paragraph{Remark on stability.}
The explicit scheme is simple and pedagogically useful, but it may
require a small time step for stability.
In practice, implicit or Crank--Nicolson schemes are often preferred.
For a first implementation, however, the explicit scheme is sufficient
to illustrate the main ideas.

\paragraph{Non-vectorized Python implementation.}
A first implementation, written in the most direct way, is as follows.

\begin{lstlisting}[language=Python]
import numpy as np

def sigma_loc(t, s, sigma0=0.2, alpha=0.3, K=100.0):
    # Example of a toy local volatility surface
    return sigma0 * (1.0 + alpha * (s - K) / K)

def price_pde_local_vol_non_vectorized(
    K=100.0, T=1.0, r=0.02, Smax=300.0,
    Nt=1000, Ns=400
):
    dt = T / Nt
    ds = Smax / Ns

    t_grid = np.linspace(0.0, T, Nt + 1)
    s_grid = np.linspace(0.0, Smax, Ns + 1)

    U = np.zeros((Nt + 1, Ns + 1))

    # terminal condition
    for i in range(Ns + 1):
        U[Nt, i] = max(s_grid[i] - K, 0.0)

    # backward induction
    for n in range(Nt - 1, -1, -1):
        t = t_grid[n]

        # boundary conditions
        U[n, 0] = 0.0
        U[n, Ns] = Smax - K * np.exp(-r * (T - t))

        for i in range(1, Ns):
            s = s_grid[i]
            sig = sigma_loc(t, s, K=K)

            dU = (U[n+1, i+1] - U[n + 1, i - 1]) / (2.0 * ds)
            d2U = (U[n+1, i+1] - 2.0 * U[n+1,i] + U[n+1, i-1]) / (ds ** 2)

            U[n,i] = U[n+1,i] + dt * (
                r * s * dU
                + 0.5 * sig**2 * s**2 * d2U
                - r * U[n+1, i]
            )

    return t_grid, s_grid, U
\end{lstlisting}

\paragraph{Vectorized Python implementation.}
The same scheme can be written in a more efficient vectorized form.

\begin{lstlisting}[language=Python]
import numpy as np

def sigma_loc_vec(t, s, sigma0=0.2, alpha=0.3, K=100.0):
    return sigma0 * (1.0 + alpha * (s - K) / K)

def price_pde_local_vol_vectorized(
    K=100.0, T=1.0, r=0.02, Smax=300.0,
    Nt=1000, Ns=400
):
    dt = T / Nt
    ds = Smax / Ns

    t_grid = np.linspace(0.0, T, Nt + 1)
    s_grid = np.linspace(0.0, Smax, Ns + 1)

    U = np.zeros((Nt + 1, Ns + 1))
    U[Nt, :] = np.maximum(s_grid - K, 0.0)

    s_inner = s_grid[1:Ns]

    for n in range(Nt - 1, -1, -1):
        t = t_grid[n]

        U[n, 0] = 0.0
        U[n, Ns] = Smax - K * np.exp(-r * (T - t))

        sig = sigma_loc_vec(t, s_inner, K=K)

        dU = (U[n + 1, 2:Ns + 1] - U[n + 1, 0:Ns - 1]) / (2.0 * ds)
        d2U = (U[n+1,2:Ns+1]-2.0*U[n+1,1:Ns]+U[n+1,0:Ns-1])/(ds**2)

        U[n, 1:Ns] = U[n + 1, 1:Ns] + dt * (
            r * s_inner * dU
            + 0.5 * sig**2 * s_inner**2 * d2U
            - r * U[n + 1, 1:Ns]
        )

    return t_grid, s_grid, U
\end{lstlisting}

\subsubsection{Numerical computation of the delta}

Once the price surface has been computed on the grid, we approximate the
delta by finite differences in the stock variable.

\paragraph{Central difference approximation.}
For interior points \(1\le i\le N_s-1\), we use
\[
\Delta(t_n,s_i)
\approx
\frac{u_{i+1}^n-u_{i-1}^n}{2\Delta s}.
\]

Thus the delta is obtained directly from the discrete price surface.

\paragraph{Boundary treatment.}
Near the boundaries, one may use one-sided finite differences, for
instance
\[
\Delta(t_n,s_0)
\approx
\frac{u_1^n-u_0^n}{\Delta s},
\qquad
\Delta(t_n,s_{N_s})
\approx
\frac{u_{N_s}^n-u_{N_s-1}^n}{\Delta s}.
\]

\paragraph{Numerical delta on the whole grid.}
The full discrete delta surface can be computed from the price matrix
\(U\).

\paragraph{Non-vectorized Python implementation.}
\begin{lstlisting}[language=Python]
def compute_delta_grid_non_vectorized(U, ds):
    Nt, Ns_plus_one = U.shape
    Ns = Ns_plus_one - 1

    Delta = np.zeros_like(U)

    for n in range(Nt):
        Delta[n, 0] = (U[n, 1] - U[n, 0]) / ds
        Delta[n, Ns] = (U[n, Ns] - U[n, Ns - 1]) / ds

        for i in range(1, Ns):
            Delta[n, i] = (U[n, i + 1] - U[n, i - 1]) / (2.0 * ds)

    return Delta
\end{lstlisting}

\paragraph{Vectorized Python implementation.}
\begin{lstlisting}[language=Python]
def compute_delta_grid_vectorized(U, ds):
    Delta = np.zeros_like(U)
    Delta[:, 0] = (U[:, 1] - U[:, 0]) / ds
    Delta[:, -1] = (U[:, -1] - U[:, -2]) / ds
    Delta[:, 1:-1] = (U[:, 2:] - U[:, :-2]) / (2.0 * ds)
    return Delta
\end{lstlisting}

\paragraph{Interpolating the delta along a simulated path.}
In hedging applications, the stock value \(S_t\) will typically not lie
exactly on the spatial grid.
We therefore need an interpolation rule.
The simplest choice is linear interpolation between neighboring grid
points.

\paragraph{Python helper function.}
\begin{lstlisting}[language=Python]
def interp_1d(x_grid, y_grid, x):
    return np.interp(x, x_grid, y_grid)
\end{lstlisting}

Thus, if \(S_t=s\), we approximate
\[
\Delta(t,s)
\]
by linearly interpolating in the row of the delta grid corresponding to
time \(t\).

\subsubsection{Discrete-time rebalancing of the hedging portfolio}

We now explain how the hedging strategy is implemented in discrete time.

\paragraph{Discrete-time setting.}
Fix a trading grid
\[
0=t_0<t_1<\cdots<t_N=T.
\]
At each rebalancing date \(t_n\), we compute a hedge ratio
\[
\Delta_{t_n},
\]
and adjust the portfolio holdings accordingly.

\paragraph{Portfolio decomposition.}
Suppose that we are short one option and hedge with:
\begin{itemize}[leftmargin=1.4em]
    \item \(\Delta_{t_n}\) shares of stock,
    \item a cash position \(B_{t_n}^{\mathrm{cash}}\).
\end{itemize}

The hedging portfolio value is
\[
\Pi_{t_n}
=
\Delta_{t_n}S_{t_n}
+
B_{t_n}^{\mathrm{cash}}.
\]

At inception, if the option is sold for \(u(0,S_0)\), we typically start
with
\[
\Pi_0=u(0,S_0),
\]
and choose
\[
B_0^{\mathrm{cash}}=u(0,S_0)-\Delta_0 S_0.
\]

\paragraph{Portfolio update between two dates.}
Between \(t_n\) and \(t_{n+1}\), the stock position is kept constant,
and the cash account accrues interest at rate \(r\).
Thus
\[
\Pi_{t_{n+1}}^{-}
=
e^{r\Delta t}
\big(\Pi_{t_n}-\Delta_{t_n}S_{t_n}\big)
+
\Delta_{t_n}S_{t_{n+1}},
\]
where
\[
\Delta t=t_{n+1}-t_n.
\]

The superscript ``\(-\)'' indicates the portfolio value just before
rebalancing at time \(t_{n+1}\).

\paragraph{Rebalancing.}
At time \(t_{n+1}\), we compute the new hedge ratio
\[
\Delta_{t_{n+1}},
\]
and we update the cash account so that the new portfolio decomposition is
\[
\Pi_{t_{n+1}}
=
\Delta_{t_{n+1}}S_{t_{n+1}}
+
B_{t_{n+1}}^{\mathrm{cash}}.
\]

In other words, the total portfolio value is unchanged by rebalancing,
but its decomposition between stock and cash is modified.

\paragraph{Non-vectorized Python implementation.}
The following function simulates one discretely rebalanced hedging
portfolio along one stock path, given a function that returns the hedge
ratio.

\begin{lstlisting}[language=Python]
def hedge_path_non_vectorized(
    S_path, t_grid, option_price_0, delta_func, r
):
    N = len(t_grid) - 1
    dt = t_grid[1] - t_grid[0]

    delta = delta_func(t_grid[0], S_path[0])
    cash = option_price_0 - delta * S_path[0]

    for n in range(N):
        # portfolio evolves from t_n to t_{n+1}
        cash = cash * np.exp(r * dt)
        portfolio = delta * S_path[n + 1] + cash

        # rebalance at t_{n+1}
        delta = delta_func(t_grid[n + 1], S_path[n + 1])
        cash = portfolio - delta * S_path[n + 1]

    return portfolio
\end{lstlisting}

\paragraph{Vectorized idea.}
For a large number of simulated paths, one can vectorize the entire
update by working with arrays of stock values, deltas, and cash
positions.
This is essential in Monte Carlo studies of hedging P\&L.

\begin{lstlisting}[language=Python]
def hedge_paths_vectorized(
    S_paths, t_grid, option_price_0, delta_func, r
):
    M, N_plus_one = S_paths.shape
    N = N_plus_one - 1
    dt = t_grid[1] - t_grid[0]

    delta = delta_func(t_grid[0], S_paths[:, 0])
    cash = option_price_0 - delta * S_paths[:, 0]

    for n in range(N):
        cash = cash * np.exp(r * dt)
        portfolio = delta * S_paths[:, n + 1] + cash

        delta = delta_func(t_grid[n + 1], S_paths[:, n + 1])
        cash = portfolio - delta * S_paths[:, n + 1]

    return portfolio
\end{lstlisting}

\paragraph{Black--Scholes versus local volatility implementation.}
The only difference between the two hedging strategies lies in the
function \texttt{delta\_func}:
\begin{itemize}[leftmargin=1.4em]
    \item for Black--Scholes, it is given by the explicit formula
    \(N(d_1)\),
    \item for local volatility, it is obtained numerically by
    interpolation from the delta grid computed from the PDE.
\end{itemize}

\subsection{Model risk: what if the true dynamics are stochastic volatility?}

We now introduce a different model, not for hedging, but for generating
the \emph{true} stock dynamics in our numerical experiment.

The purpose is to compare two hedging strategies:
\begin{itemize}[leftmargin=1.4em]
    \item a delta hedge computed under Black--Scholes,
    \item a delta hedge computed under local volatility,
\end{itemize}
when neither of these models is actually the true one.

This is a simple and concrete way to illustrate model risk.

\subsubsection{A reference stochastic volatility model}

As a reference ``true'' dynamics, we consider a Heston-type stochastic
volatility model:
\[
\mathrm{d}S_t
=
rS_t\,\mathrm{d}t
+
\sqrt{V_t}\,S_t\,\mathrm{d}W_t^{(1)},
\]
\[
\mathrm{d}V_t
=
\kappa(\theta-V_t)\,\mathrm{d}t
+
\eta\sqrt{V_t}\,\mathrm{d}W_t^{(2)},
\]
with
\[
\mathrm{d}\langle W^{(1)},W^{(2)}\rangle_t=\rho\,\mathrm{d}t.
\]

\paragraph{Interpretation.}
In this model:
\begin{itemize}[leftmargin=1.4em]
    \item the stock still has drift \(r\) under the risk-neutral measure,
    \item but the instantaneous variance \(V_t\) is itself random,
    \item and the correlation parameter \(\rho\) creates asymmetry between
    spot and volatility shocks.
\end{itemize}

\paragraph{Why use such a model for the experiment?}
This is a natural benchmark for model risk because:
\begin{itemize}[leftmargin=1.4em]
    \item Black--Scholes is too simple compared with it,
    \item local volatility may capture the current smile but not the full
    stochastic dynamics of volatility,
    \item and the resulting hedging experiment is rich enough to produce
    meaningful P\&L differences.
\end{itemize}

\subsubsection{Simulating the true dynamics}

To compare hedging strategies, we simulate stock paths under the
stochastic volatility model.

\paragraph{Euler-type scheme.}
On a time grid \(t_n=n\Delta t\), we may use the discretization
\[
S_{n+1}
=
S_n
+
rS_n\Delta t
+
\sqrt{V_n^+}\,S_n\,\Delta W_n^{(1)},
\]
\[
V_{n+1}
=
V_n
+
\kappa(\theta-V_n)\Delta t
+
\eta\sqrt{V_n^+}\,\Delta W_n^{(2)},
\]
where
\[
V_n^+=\max(V_n,0).
\]

To generate correlated Brownian increments, we may write
\[
\Delta W_n^{(1)}=\sqrt{\Delta t}\,Z_n^{(1)},
\]
\[
\Delta W_n^{(2)}
=
\sqrt{\Delta t}\big(
\rho Z_n^{(1)}+\sqrt{1-\rho^2}\,Z_n^{(2)}
\big),
\]
where \(Z_n^{(1)}\) and \(Z_n^{(2)}\) are independent standard normal
random variables.

\paragraph{Non-vectorized Python implementation.}
\begin{lstlisting}[language=Python]
def simulate_heston_path_non_vectorized(
    S0, V0, r, kappa, theta, eta, rho, T, N
):
    dt = T / N
    t_grid = np.linspace(0.0, T, N + 1)

    S = np.zeros(N + 1)
    V = np.zeros(N + 1)

    S[0] = S0
    V[0] = V0

    for n in range(N):
        z1 = np.random.randn()
        z2 = np.random.randn()

        dW1 = np.sqrt(dt) * z1
        dW2 = np.sqrt(dt) * (rho * z1 + np.sqrt(1.0 - rho**2) * z2)

        V_pos = max(V[n], 0.0)

        S[n + 1] = S[n] + r * S[n] * dt + np.sqrt(V_pos) * S[n] * dW1
        V[n+1]=V[n]+kappa*(theta-V[n])*dt+eta*np.sqrt(V_pos)*dW2

    return t_grid, S, V
\end{lstlisting}

\paragraph{Vectorized Python implementation.}
\begin{lstlisting}[language=Python]
def simulate_heston_paths_vectorized(
    M, S0, V0, r, kappa, theta, eta, rho, T, N
):
    dt = T / N
    t_grid = np.linspace(0.0, T, N + 1)

    S = np.zeros((M, N + 1))
    V = np.zeros((M, N + 1))

    S[:, 0] = S0
    V[:, 0] = V0

    for n in range(N):
        z1 = np.random.randn(M)
        z2 = np.random.randn(M)

        dW1 = np.sqrt(dt) * z1
        dW2 = np.sqrt(dt) * (rho * z1 + np.sqrt(1.0 - rho**2) * z2)

        V_pos = np.maximum(V[:, n], 0.0)

        S[:, n + 1] = S[:, n] + r * S[:, n] * dt + np.sqrt(V_pos) * S[:, n] * dW1
        V[:,n+1]=V[:,n]+kappa*(theta-V[:,n])*dt+eta*np.sqrt(V_pos)*dW2
    return t_grid, S, V
\end{lstlisting}

\subsubsection{Hedging with the Black--Scholes delta}

The first hedging strategy uses the Black--Scholes model as if it were
the true one.

\paragraph{Hedge ratio.}
At each rebalancing date \(t_n\), the hedge ratio is
\[
\Delta_{t_n}^{BS}
=
N(d_1(t_n,S_{t_n})),
\]
where \(d_1\) is computed using the Black--Scholes volatility parameter.

\paragraph{Interpretation.}
This strategy is simple and explicit, but it ignores:
\begin{itemize}[leftmargin=1.4em]
    \item the market smile,
    \item the state dependence of local volatility,
    \item and the stochastic fluctuations of the true variance process.
\end{itemize}

Thus it is best understood as a benchmark hedge.

\paragraph{Python implementation.}
\begin{lstlisting}[language=Python]
from math import log, sqrt, exp
from scipy.stats import norm

def bs_delta(t, s, K, T, r, sigma):
    tau = max(T - t, 1e-12)
    d1 = (np.log(s / K) + (r + 0.5 * sigma**2) * tau) / (sigma * np.sqrt(tau))
    return norm.cdf(d1)
\end{lstlisting}

\subsubsection{Hedging with the local volatility delta}

The second hedging strategy uses the local volatility model calibrated
from the market vanilla surface.

\paragraph{Hedge ratio.}
At each date \(t_n\), we compute
\[
\Delta_{t_n}^{LV}
=
\partial_s u^{LV}(t_n,S_{t_n}),
\]
where \(u^{LV}\) solves the local volatility PDE.

Numerically, this means:
\begin{itemize}[leftmargin=1.4em]
    \item solve the PDE on a grid,
    \item compute the discrete delta grid,
    \item interpolate the delta at the realized stock level.
\end{itemize}

\paragraph{Interpretation.}
This strategy is more market-consistent than the Black--Scholes hedge
because it incorporates the current implied-volatility smile through the
local volatility surface.

However, it still assumes that the true dynamics are given by a local
volatility model, whereas in our experiment the true dynamics are
stochastic volatility.

\paragraph{Python wrapper.}
A simple implementation is to use a function that interpolates the delta
surface at a given time and spot.

\begin{lstlisting}[language=Python]
def build_local_vol_delta_function(t_grid, s_grid, Delta_grid):
    def delta_func(t, s):
        # time interpolation: nearest grid point for simplicity
        n = np.searchsorted(t_grid, t)
        n = min(max(n, 0), len(t_grid) - 1)
        return np.interp(s, s_grid, Delta_grid[n, :])
    return delta_func
\end{lstlisting}

A more refined implementation would interpolate both in time and in
space, but the idea is already clear with this simple version.

\subsection{Comparing hedging P\&L}

We now explain how to compare the performance of the two hedging
strategies.

\subsubsection{Definition of the hedging P\&L}

Suppose that we sell one European call at time \(0\) for the model price
\[
u(0,S_0),
\]
and then hedge dynamically up to maturity.

Let \(\Pi_T\) denote the terminal value of the hedging portfolio.
The terminal payoff of the short option position is
\[
-(S_T-K)^+.
\]

Thus the final hedging profit and loss is
\[
\boxed{
\mathrm{PnL}
=
\Pi_T-(S_T-K)^+.
}
\]

\paragraph{Interpretation.}
If the hedge were perfect and rebalancing continuous, then in the ideal
model one would have
\[
\mathrm{PnL}=0.
\]
In practice, however, several sources of error remain:
\begin{itemize}[leftmargin=1.4em]
    \item discrete rebalancing,
    \item numerical approximation,
    \item and, most importantly here, model misspecification.
\end{itemize}

\paragraph{Two P\&L distributions.}
We will therefore obtain:
\begin{itemize}[leftmargin=1.4em]
    \item one distribution of P\&L from the Black--Scholes hedge,
    \item one distribution of P\&L from the local-volatility hedge.
\end{itemize}

The comparison of these two distributions is the main output of the
numerical experiment.

\subsubsection{Pathwise comparison of the two strategies}

For each simulated path of the true stochastic volatility model, we can:
\begin{itemize}[leftmargin=1.4em]
    \item compute the terminal payoff,
    \item compute the terminal value of the Black--Scholes hedging
    portfolio,
    \item compute the terminal value of the local-volatility hedging
    portfolio,
    \item and compare the two resulting P\&Ls.
\end{itemize}

\paragraph{Why pathwise comparison is useful.}
Pathwise comparison makes the model-risk issue very concrete.
Even for the same realized stock path, the two hedging strategies may
perform differently because the hedge ratios used along the path are not
the same.

\paragraph{Typical numerical outputs.}
Useful pathwise diagnostics include:
\begin{itemize}[leftmargin=1.4em]
    \item plotting one or several stock paths,
    \item plotting the corresponding hedge ratios through time,
    \item plotting the cumulative value of the hedging portfolios,
    \item and comparing the final pathwise errors.
\end{itemize}

\paragraph{Illustrative Python snippet.}
\begin{lstlisting}[language=Python]
payoff = np.maximum(S_paths[:, -1] - K, 0.0)

pnl_bs = portfolio_bs - payoff
pnl_lv = portfolio_lv - payoff
\end{lstlisting}

\subsubsection{Distribution of the hedging error}

Once a large number of paths has been simulated, we can study the
distribution of the hedging error for each strategy.

\paragraph{Quantities of interest.}
Typical summary statistics include:
\begin{itemize}[leftmargin=1.4em]
    \item the mean P\&L,
    \item the standard deviation,
    \item quantiles,
    \item the probability of large losses.
\end{itemize}

\paragraph{Visual comparison.}
A natural way to compare the two strategies is to display:
\begin{itemize}[leftmargin=1.4em]
    \item histograms of the two P\&L distributions,
    \item boxplots,
    \item or cumulative distribution curves.
\end{itemize}

\paragraph{Python example.}
\begin{lstlisting}[language=Python]
print("BS hedge:")
print("mean =", np.mean(pnl_bs))
print("std  =", np.std(pnl_bs))
print("q05  =", np.quantile(pnl_bs, 0.05))
print("q95  =", np.quantile(pnl_bs, 0.95))

print("Local vol hedge:")
print("mean =", np.mean(pnl_lv))
print("std  =", np.std(pnl_lv))
print("q05  =", np.quantile(pnl_lv, 0.05))
print("q95  =", np.quantile(pnl_lv, 0.95))
\end{lstlisting}

One may then add:
\begin{lstlisting}[language=Python]
import matplotlib.pyplot as plt

plt.hist(pnl_bs, bins=60, alpha=0.5, label="BS hedge")
plt.hist(pnl_lv, bins=60, alpha=0.5, label="Local vol hedge")
plt.legend()
plt.title("Distribution of hedging P&L")
plt.show()
\end{lstlisting}

\paragraph{What we are looking for.}
A hedging strategy is typically considered better if:
\begin{itemize}[leftmargin=1.4em]
    \item its P\&L is more concentrated around zero,
    \item its variance is smaller,
    \item and its tail losses are less severe.
\end{itemize}

\subsubsection{Interpretation of the results}

The purpose of the experiment is not to declare one model universally
superior, but to illustrate how the choice of hedging model affects the
distribution of the hedging error.

\paragraph{Expected qualitative outcome.}
Since the true dynamics are stochastic volatility:
\begin{itemize}[leftmargin=1.4em]
    \item the Black--Scholes hedge is based on a misspecified constant
    volatility model,
    \item the local-volatility hedge is based on a model calibrated to
    the current smile, but still misspecified dynamically,
    \item so neither hedge is perfect.
\end{itemize}

However, the local-volatility hedge may be expected to perform better in
some scenarios because it incorporates more information about the market
surface than the Black--Scholes hedge.

\paragraph{Important nuance.}
This is not a theorem.
Depending on:
\begin{itemize}[leftmargin=1.4em]
    \item the option strike and maturity,
    \item the shape of the local volatility surface,
    \item the stochastic volatility parameters,
    \item and the rebalancing frequency,
\end{itemize}
the relative performance of the two hedges may vary.

\paragraph{Main lesson from the experiment.}
The key point is that:
\begin{quote}
a hedge is always computed inside a model, but its performance is judged
in the true market dynamics.
\end{quote}

Thus pricing models should not be evaluated only through the price they
produce, but also through the quality of the hedging strategies they
generate.

\subsection{Main lessons from the European option example}

This first example contains most of the central ideas of option pricing
and hedging.

\paragraph{First lesson: pricing and hedging are linked.}
Once a model for the underlying is specified, one obtains:
\begin{itemize}[leftmargin=1.4em]
    \item a risk-neutral pricing formula,
    \item a pricing PDE,
    \item and a delta hedge.
\end{itemize}

Thus a pricing model is not only a machine for generating a price: it is
also a machine for generating a hedging strategy.

\paragraph{Second lesson: the same contract can lead to different hedges.}
We considered the same European call option under two models:
\begin{itemize}[leftmargin=1.4em]
    \item Black--Scholes,
    \item local volatility.
\end{itemize}

The payoff did not change, but the price function changed, and therefore
the delta changed as well.

Thus the hedge ratio is fundamentally model-dependent.

\paragraph{Third lesson: market consistency matters.}
The Black--Scholes hedge is simple, explicit, and conceptually central.
The local-volatility hedge is more complex, but it is calibrated to the
current market smile.

Hence the local-volatility model can be seen as a market-consistent
extension of the Black--Scholes benchmark.

\paragraph{Fourth lesson: model risk remains even for a vanilla product.}
Even for a simple European call, if the true underlying dynamics are
stochastic volatility, neither the Black--Scholes hedge nor the
local-volatility hedge is exact.

Therefore, a model should be judged not only by its pricing formula, but
also by the distribution of the hedging error it produces in realistic
market scenarios.

\paragraph{Final transition.}
This concludes the first case study.
It already illustrates the main relation between:
\[
\text{model}
\longrightarrow
\text{price}
\longrightarrow
\text{PDE}
\longrightarrow
\text{delta}
\longrightarrow
\text{hedging P\&L}.
\]

In the next example, we will move to a more path-dependent contract, for
which the choice of model becomes even more important.

% ============================================================
\section{Path-dependent option: the lookback option}
% ============================================================

\subsection{The contract and why path dependence matters}

We now move from a European vanilla option to a path-dependent option.
The purpose of this second example is to understand how the pricing and
hedging problem changes when the payoff depends not only on the terminal
stock price \(S_T\), but also on the past trajectory of the stock.

\paragraph{The contract.}
We consider a fixed-strike lookback call option with maturity \(T\) and
strike \(K\).
Its payoff is
\[
(M_T-K)^+,
\qquad
\text{where }
M_T:=\max_{0\le u\le T} S_u.
\]

Thus the holder receives the positive part of the difference between the
running maximum of the stock over the whole life of the contract and the
strike \(K\).

\paragraph{Why this product is different from a vanilla European call.}
For a European call, the payoff is
\[
(S_T-K)^+,
\]
so only the terminal value \(S_T\) matters.

For the lookback option, this is no longer true.
The payoff depends on the whole path of the stock through the quantity
\[
M_T=\max_{0\le u\le T}S_u.
\]

Therefore, two stock paths with the same terminal value \(S_T\) may lead
to different option payoffs if they have reached different maxima before
maturity.

\paragraph{A simple illustration.}
Suppose that two stock paths both end at the same terminal level
\[
S_T=100.
\]
If the first path has reached a maximum of \(120\) during the life of
the option, while the second path has never gone above \(105\), then the
lookback option will produce two different payoffs:
\[
(120-K)^+
\qquad\text{and}\qquad
(105-K)^+.
\]

This illustrates the central feature of path dependence:
the current spot price alone does not contain enough information to
price the contract.

\paragraph{Why path dependence matters for pricing.}
For a European option in a Markovian model, the price can often be
written as a function of the current time and current stock price only:
\[
u(t,S_t).
\]
For a lookback option, this is no longer sufficient.
At time \(t\), the future payoff depends not only on the current stock
price \(S_t\), but also on the current running maximum
\[
M_t:=\max_{0\le u\le t}S_u.
\]

Hence the natural state variable is no longer one-dimensional.
To price the contract, we must keep track of both:
\begin{itemize}[leftmargin=1.4em]
    \item the current stock price \(S_t\),
    \item and the current running maximum \(M_t\).
\end{itemize}

Thus the option price will naturally be written as
\[
u(t,S_t,M_t).
\]

\paragraph{Why path dependence matters for hedging.}
The same issue appears in the hedging problem.
For a European call, the hedge ratio depends on the current stock price
and time.
For a lookback option, the hedge ratio is also affected by the position
of the current stock relative to the running maximum.

Intuitively:
\begin{itemize}[leftmargin=1.4em]
    \item if the current stock is far below the running maximum, the
    chance of creating a new maximum before maturity may be small;
    \item if the stock is close to the running maximum, then a small
    upward move may change the future payoff structure significantly.
\end{itemize}

Thus the delta of a lookback option is usually more sensitive and more
unstable than the delta of a vanilla European call.

\paragraph{Main objective of this example.}
In this second example, we want to understand how the general
pricing--hedging methodology changes when the payoff becomes
path-dependent.

As in the European-option case, we will follow the same sequence:
\begin{itemize}[leftmargin=1.4em]
    \item define the contract,
    \item choose a pricing model,
    \item write the risk-neutral pricing formula,
    \item derive the PDE in the relevant state variables,
    \item discuss the numerical implementation,
    \item and compare hedging strategies through their final P\&L.
\end{itemize}

The main difference is that the state space and the hedging problem are
now richer.

\subsection{Choice of the model}

A path-dependent option such as a lookback option raises immediately the
question of model choice.

As before, we could start from a simple benchmark model and then move to
a richer one.
However, the role of the model is even more important here than in the
European case, because the payoff depends on the whole path of the
underlying.

We now explain how to think about this choice.

\subsubsection{Why a simple Black--Scholes benchmark is still useful}

Even though the lookback option is more complex than a vanilla option,
the Black--Scholes model remains a very useful benchmark.

\paragraph{Black--Scholes benchmark.}
Under the risk-neutral measure, the Black--Scholes model assumes
\[
\mathrm{d}S_t=rS_t\,\mathrm{d}t+\sigma S_t\,\mathrm{d}W_t,
\]
with constant volatility \(\sigma\).

This model is attractive as a first benchmark for several reasons.

\paragraph{First reason: conceptual simplicity.}
The model is familiar from the European-option case.
Students already know:
\begin{itemize}[leftmargin=1.4em]
    \item the risk-neutral dynamics,
    \item the pricing logic,
    \item the PDE derivation,
    \item and the meaning of dynamic hedging.
\end{itemize}

Thus using Black--Scholes first allows us to focus on what is genuinely
new, namely the path dependence, without changing too many ingredients
at once.

\paragraph{Second reason: the new difficulty comes from the payoff, not yet from the model.}
Even under Black--Scholes, the lookback option is already a much richer
problem than the European call, because one must keep track of the
running maximum.
Thus the first source of difficulty is already present before any
volatility smile or stochastic volatility effect is introduced.

\paragraph{Third reason: it provides a clean benchmark for comparison.}
Later, when we introduce a richer model, we will want to compare:
\begin{itemize}[leftmargin=1.4em]
    \item prices,
    \item hedging strategies,
    \item and hedging P\&L.
\end{itemize}

For this, it is essential to have a simple reference model.
Black--Scholes plays that role.

\paragraph{Fourth reason: numerical implementation remains manageable.}
In the Black--Scholes setting, the dynamics are simple and simulation is
straightforward.
Thus the additional complexity of the product can be handled clearly:
\begin{itemize}[leftmargin=1.4em]
    \item in Monte Carlo, one must simulate the path and track the
    running maximum;
    \item in PDE form, one must enlarge the state space to include the
    running maximum.
\end{itemize}

So Black--Scholes is still pedagogically very useful.

\paragraph{Important limitation.}
Of course, Black--Scholes is only a benchmark.
It assumes constant volatility and therefore ignores smile, skew, and
the possible randomness of future volatility.
For a path-dependent product, this can be a serious limitation.

That is exactly why we need to think more carefully about richer models.

\subsubsection{Why path dependence makes the problem more complex}

The main source of additional complexity comes from the structure of the
payoff itself.

\paragraph{Loss of a one-dimensional state description.}
For a European option in a Markovian stock model, the option price can
be written as
\[
u(t,s),
\]
because the current stock price \(s\) contains all the information
needed about the future distribution of the payoff.

For a lookback option, this is no longer true.
At time \(t\), the future payoff depends on both:
\begin{itemize}[leftmargin=1.4em]
    \item the current stock price \(S_t\),
    \item and the current running maximum
    \[
    M_t=\max_{0\le u\le t}S_u.
    \]
\end{itemize}

Thus the problem is no longer one-dimensional.
The natural price function becomes
\[
u(t,s,m),
\qquad \text{with } s\le m.
\]

\paragraph{The product is path-sensitive even when the model is simple.}
This is a very important point:
even if the stock follows the simplest possible diffusion, the option is
still nontrivial because its payoff depends on the path.

So the complexity here does not come only from the model.
It already comes from the contract.

\paragraph{The running maximum creates a moving boundary effect.}
The running maximum does not evolve like an ordinary diffusion variable.
Its dynamics are constrained:
\begin{itemize}[leftmargin=1.4em]
    \item if \(S_t<M_t\), then the running maximum stays constant;
    \item if \(S_t=M_t\) and the stock moves upward, then the running
    maximum also moves upward.
\end{itemize}

Thus the additional state variable \(M_t\) is not just another
independent factor.
It is a path functional tied to the stock process.

This creates additional subtleties in PDE formulations and in numerical
methods.

\paragraph{The hedge is more sensitive to the path history.}
For a vanilla call, the delta depends mainly on the current spot and
time to maturity.
For a lookback option, the delta depends also on how close the stock is
to the running maximum.

In particular, if \(S_t\) is very close to \(M_t\), then a small upward
move may create a new maximum, which changes the future payoff in a
highly nonlinear way.
This can make the hedge more fragile and more sensitive to model choice.

\paragraph{Discrete hedging becomes more delicate.}
Since the payoff depends on the running maximum, a continuous-time model
may react very differently from a discrete-time hedging strategy.
Missing a short-lived move above a previous maximum can have a strong
effect on the final payoff.
Thus discretization error is often more important for path-dependent
products than for standard European options.

\paragraph{Main message.}
Even before discussing local volatility or stochastic volatility, the
lookback option is already more difficult than the European call because
the relevant state variable is no longer just the current stock level.
Path dependence forces us to enlarge the state space and to rethink both
pricing and hedging.

\subsubsection{Why a richer volatility model may matter more for this product}

The next question is whether a simple Black--Scholes benchmark is enough
in practice for a path-dependent product such as a lookback option.

The short answer is: often, no.

\paragraph{Sensitivity to the whole path.}
A European option depends only on the terminal marginal distribution of
\(S_T\).
A lookback option depends on the whole trajectory of \(S\), because the
payoff involves the maximum over time.

This means that the model must describe not only where the stock ends,
but also how it gets there.

\paragraph{Why volatility modeling matters more here.}
Different volatility models may generate similar prices for some liquid
European options at time \(0\), but very different path behavior for the
stock.
In particular, the probability that the stock reaches high levels before
maturity can be strongly affected by:
\begin{itemize}[leftmargin=1.4em]
    \item the shape of the implied volatility smile,
    \item the local behavior of volatility as a function of spot,
    \item or the stochastic evolution of volatility itself.
\end{itemize}

Since a lookback option is directly sensitive to extreme values reached
by the stock, it is naturally more model-sensitive than a vanilla call.

\paragraph{Why local volatility is a natural next step.}
A local volatility model has the form
\[
\mathrm{d}S_t
=
rS_t\,\mathrm{d}t
+
\sigma_{\mathrm{loc}}(t,S_t)S_t\,\mathrm{d}W_t.
\]

Compared with Black--Scholes, this keeps the model one-dimensional and
Markovian in the stock, but it incorporates the market smile through the
state-dependent function \(\sigma_{\mathrm{loc}}(t,s)\).

For a path-dependent product, this is attractive because:
\begin{itemize}[leftmargin=1.4em]
    \item the model remains relatively tractable,
    \item the stock dynamics are consistent with the current vanilla
    surface,
    \item and the probability of hitting high levels can change
    significantly compared with the constant-volatility case.
\end{itemize}

Thus local volatility is a natural market-consistent extension of the
Black--Scholes benchmark.

\paragraph{Why even local volatility may still be insufficient.}
At the same time, local volatility remains a deterministic volatility
model.
It fits today's smile, but it does not introduce an independent random
volatility factor.

For a path-dependent product, this can matter even more than for a
vanilla option, because the future path of the stock may be strongly
affected by the future path of volatility.
This is why stochastic volatility models are often considered more
realistic for exotic products.

\paragraph{Pedagogical strategy for this section.}
For the purposes of the course, the right approach is therefore:
\begin{itemize}[leftmargin=1.4em]
    \item start from Black--Scholes as a clean benchmark,
    \item then explain why the path-dependent nature of the product makes
    model choice more important,
    \item and use a richer model, such as local volatility, to produce a
    more market-consistent pricing and hedging framework.
\end{itemize}

\paragraph{Main lesson.}
For a path-dependent product such as a lookback option, the choice of
model matters more than for a standard European call.
The reason is simple:
\begin{quote}
the contract depends on the whole path of the stock, so the model must
describe path behavior well, not only terminal distributions.
\end{quote}

\subsection{State variables and pricing problem}

We now explain how the pricing problem must be reformulated once the
payoff depends on the running extremum of the stock price.

For a standard European option, the current stock price \(S_t\) is
enough to describe the state of the problem.
For a lookback option, this is no longer true.
The relevant state variable must contain additional information about
the past path of the stock.

\subsubsection{Running maximum / running minimum as additional state variables}

The key observation is that the payoff of a lookback option depends on a
running extremum of the stock process.

For a fixed-strike lookback call, the payoff is
\[
(M_T-K)^+,
\qquad
M_T:=\max_{0\le u\le T}S_u.
\]
So the natural additional state variable is the running maximum
\[
M_t:=\max_{0\le u\le t}S_u.
\]

Similarly, for a lookback put or for other path-dependent contracts, one
may need the running minimum
\[
m_t:=\min_{0\le u\le t}S_u.
\]

\paragraph{Why this additional variable is necessary.}
If at time \(t\) we know only the current stock value \(S_t\), then in
general we do not know the current running maximum \(M_t\).
But the future payoff depends on it.

Indeed, two trajectories may share the same current value \(S_t\) while
having reached different historical maxima.
Since the future payoff depends on the maximum over the whole interval
\([0,T]\), these two states are not equivalent from the point of view of
pricing.

\paragraph{State constraints.}
By construction, we always have
\[
S_t\le M_t.
\]
Thus the relevant state space is not the whole plane, but the region
\[
\{(s,m)\in\R_+^2:\ s\le m\}.
\]

If we were instead working with a running minimum, we would have
\[
m_t\le S_t.
\]

\paragraph{Interpretation.}
The running maximum \(M_t\) summarizes the part of the past trajectory
that is relevant for future pricing.
It is not an independent stochastic factor like a second Brownian
motion; rather, it is a path functional generated by the stock process
itself.

This is precisely what makes the problem path-dependent but still
amenable to a Markovian reformulation.

\subsubsection{Markovian reformulation of the problem}

Even though the original payoff depends on the whole past path of the
stock, one can recover a Markovian structure by enlarging the state
space.

\paragraph{From a path-dependent problem to a Markovian one.}
If the stock alone follows a Markov process under the pricing model, then
the pair
\[
(S_t,M_t)
\]
is itself Markov.
In other words, once we know:
\begin{itemize}[leftmargin=1.4em]
    \item the current stock level \(S_t\),
    \item and the current running maximum \(M_t\),
\end{itemize}
we have all the information needed to price the option from time \(t\)
onward.

Thus the option price can be written as a deterministic function
\[
u(t,s,m),
\qquad \text{with } s\le m.
\]

\paragraph{Why this is useful.}
This enlarged Markovian formulation is the key step that allows us to:
\begin{itemize}[leftmargin=1.4em]
    \item write a risk-neutral pricing formula,
    \item derive a PDE,
    \item and implement numerical methods.
\end{itemize}

Without this state augmentation, the problem would remain truly
path-dependent and much harder to analyze through PDE methods.

\paragraph{Dynamics of the running maximum.}
The variable \(M_t\) evolves in a constrained way:
\begin{itemize}[leftmargin=1.4em]
    \item if \(S_t<M_t\), then the running maximum stays constant;
    \item if \(S_t=M_t\) and the stock moves upward, then the running
    maximum moves with the stock.
\end{itemize}

Thus \(M_t\) changes only when the stock hits a new historical maximum.

This feature will later appear in the PDE formulation through a boundary
condition on the diagonal \(s=m\).

\paragraph{A useful analogy.}
For a European option, the state variable is simply
\[
(t,S_t).
\]
For a lookback option, the state variable becomes
\[
(t,S_t,M_t).
\]

So the logic is unchanged, but the state space has gained one dimension.
This is the mathematical price one pays for path dependence.

\paragraph{The same idea for the running minimum.}
If the contract involved the running minimum instead of the running
maximum, then the appropriate Markovian state would be
\[
(t,S_t,m_t),
\qquad
m_t:=\min_{0\le u\le t}S_u.
\]

Thus the general principle is the following:
\begin{quote}
for a path-dependent payoff, one must identify the smallest collection
of path functionals that makes the problem Markov again.
\end{quote}

\subsubsection{Risk-neutral pricing formula}

Once the state variable has been enlarged, the pricing problem takes a
form very similar to the one used for European options.

\paragraph{Pricing function.}
Let
\[
u(t,s,m)
\]
denote the time-\(t\) price of the fixed-strike lookback call when
\[
S_t=s,
\qquad
M_t=m.
\]

\paragraph{Risk-neutral valuation.}
Under the chosen pricing model and under the risk-neutral measure
\(\Qbb\), the option price is given by
\[
\boxed{
u(t,s,m)
=
\E^\Qbb\!\left[
e^{-r(T-t)}(M_T-K)^+
\;\middle|\;
S_t=s,\ M_t=m
\right].
}
\]

This is the exact analogue of the risk-neutral pricing formula for a
European option, except that the conditioning now includes the running
maximum.

\paragraph{Interpretation.}
This formula says that, conditional on the current stock level \(s\) and
the current running maximum \(m\), the option price is the discounted
conditional expectation of the final payoff.

So the pricing principle itself has not changed:
it is still risk-neutral valuation.
What has changed is only the state variable.

\paragraph{Special role of the current maximum.}
At time \(t\), the future running maximum \(M_T\) can be written
informally as the larger of:
\begin{itemize}[leftmargin=1.4em]
    \item the maximum already achieved, namely \(M_t=m\),
    \item and the future maxima reached by the stock over \([t,T]\).
\end{itemize}

Thus the current value \(m\) acts as a memory variable summarizing the
relevant past information.

\paragraph{Comparison with the European case.}
For a European call, we had
\[
u(t,s)
=
\E^\Qbb\!\left[
e^{-r(T-t)}(S_T-K)^+
\;\middle|\;
S_t=s
\right].
\]

For the lookback call, the corresponding formula becomes
\[
u(t,s,m)
=
\E^\Qbb\!\left[
e^{-r(T-t)}(M_T-K)^+
\;\middle|\;
S_t=s,\ M_t=m
\right].
\]

Thus the structure is the same, but the dependence on the path is now
captured through the additional conditioning variable \(m\).

\paragraph{What comes next.}
Once this risk-neutral pricing formula is written, the next step is to
derive the PDE satisfied by the function \(u(t,s,m)\), together with the
appropriate terminal and boundary conditions.

This will show explicitly how the additional state variable enters the
pricing problem.

\subsection{Pricing PDE for the lookback option}

We now derive the PDE satisfied by the price of the lookback option once
the problem has been reformulated in the enlarged state space
\[
(t,S_t,M_t).
\]

The key point is that the state variable is now two-dimensional in
space:
\begin{itemize}[leftmargin=1.4em]
    \item the current stock price \(S_t\),
    \item the current running maximum \(M_t\).
\end{itemize}

This is the main mathematical difference with the European call case.

\subsubsection{PDE in the extended state space}

Let
\[
u(t,s,m)
\]
denote the time-\(t\) price of the fixed-strike lookback call when
\[
S_t=s,
\qquad
M_t=m,
\qquad
s\le m.
\]

We assume for the moment that the stock follows the Black--Scholes
dynamics under the risk-neutral measure:
\[
\mathrm{d}S_t=rS_t\,\mathrm{d}t+\sigma S_t\,\mathrm{d}W_t.
\]

\paragraph{Interior region.}
Consider first the region
\[
\{(t,s,m): 0<s<m\}.
\]
In this region, the stock is strictly below the running maximum.
Therefore, over an infinitesimal time interval, the running maximum does
not change unless the stock hits the boundary \(s=m\).

So, inside the region \(s<m\), the variable \(m\) is locally frozen and
the dynamics are driven only by the stock process.

\paragraph{It\^o formula in the interior.}
Assume that \(u\) is sufficiently smooth.
As long as \(s<m\), the process \(M_t\) remains locally constant, so
applying It\^o's formula to \(u(t,S_t,M_t)\) gives
\[
\mathrm{d}u(t,S_t,M_t)
=
\partial_t u\,\mathrm{d}t
+
\partial_s u\,\mathrm{d}S_t
+
\frac12 \partial_{ss}u\,\mathrm{d}\langle S\rangle_t.
\]

Since
\[
\mathrm{d}S_t=rS_t\,\mathrm{d}t+\sigma S_t\,\mathrm{d}W_t
\]
and
\[
\mathrm{d}\langle S\rangle_t=\sigma^2S_t^2\,\mathrm{d}t,
\]
we obtain
\[
\mathrm{d}u(t,S_t,M_t)
=
\Big(
\partial_t u
+
rs\,\partial_s u
+
\frac12 \sigma^2 s^2\,\partial_{ss}u
\Big)\mathrm{d}t
+
\sigma s\,\partial_s u\,\mathrm{d}W_t.
\]

\paragraph{Risk-neutral martingale condition.}
Under the risk-neutral measure, the discounted option price must be a
martingale.
Therefore, in the interior region \(s<m\), the drift of
\[
e^{-rt}u(t,S_t,M_t)
\]
must vanish.
This yields the PDE
\[
\boxed{
\partial_t u(t,s,m)
+
rs\,\partial_s u(t,s,m)
+
\frac12 \sigma^2 s^2\,\partial_{ss}u(t,s,m)
-r\,u(t,s,m)
=
0,
\qquad 0<s<m.
}
\]

\paragraph{Important observation.}
This equation looks almost identical to the Black--Scholes PDE for a
European option.
The difference is that:
\begin{itemize}[leftmargin=1.4em]
    \item the price depends on the extra variable \(m\),
    \item but inside the continuation region \(s<m\), there are no
    derivatives with respect to \(m\).
\end{itemize}

This is because the running maximum does not evolve continuously in the
interior.
It changes only when the stock reaches a new maximum, that is, on the
boundary \(s=m\).

\paragraph{Generalization to local volatility.}
If instead the stock follows a local volatility model
\[
\mathrm{d}S_t
=
rS_t\,\mathrm{d}t
+
\sigma_{\mathrm{loc}}(t,S_t)S_t\,\mathrm{d}W_t,
\]
then the same reasoning gives
\[
\boxed{
\partial_t u(t,s,m)
+
rs\,\partial_s u(t,s,m)
+
\frac12 \sigma_{\mathrm{loc}}(t,s)^2 s^2\,\partial_{ss}u(t,s,m)
-r\,u(t,s,m)
=
0,
\qquad 0<s<m.
}
\]

Thus, as in the European case, one simply replaces the constant
volatility \(\sigma\) by the local volatility function.

\subsubsection{Terminal condition}

At maturity \(T\), the contract pays
\[
(M_T-K)^+.
\]

Therefore the terminal condition for the pricing function is
\[
\boxed{
u(T,s,m)=(m-K)^+,
\qquad 0<s\le m.
}
\]

\paragraph{Interpretation.}
This terminal condition is fundamentally different from the one for a
European call.

For a European call, the payoff at maturity depends on the terminal
stock value:
\[
u(T,s)=(s-K)^+.
\]

For the lookback call, the payoff depends on the running maximum:
\[
u(T,s,m)=(m-K)^+.
\]

So at maturity, the current stock price \(s\) is not enough to determine
the payoff.
The relevant variable is the running maximum \(m\).

\paragraph{A useful remark.}
At time \(T\), one always has
\[
s=S_T\le M_T=m.
\]
Thus the payoff is determined by the largest value reached by the stock
up to maturity, not necessarily by its final position.

This is exactly the feature that makes the lookback option path
dependent.

\subsubsection{Boundary conditions}

To complete the PDE formulation, we must specify the appropriate
boundary conditions.

There are two main boundaries to consider:
\begin{itemize}[leftmargin=1.4em]
    \item the lower boundary in the stock variable,
    \item the diagonal boundary \(s=m\), where the stock is at its
    running maximum.
\end{itemize}

\paragraph{Boundary at \(s=0\).}
If the stock reaches zero, then under the Black--Scholes model it
remains at zero afterward.
The running maximum, however, remains equal to its past value \(m\).

Therefore, if \(s=0\), the value of the contract is simply the discounted
expected payoff based on the already achieved maximum:
\[
u(t,0,m)=e^{-r(T-t)}(m-K)^+.
\]

So the boundary condition at \(s=0\) is
\[
\boxed{
u(t,0,m)=e^{-r(T-t)}(m-K)^+.
}
\]

\paragraph{Interpretation.}
Once the stock is at zero, no new maximum can be reached.
The future payoff is therefore already determined by the current running
maximum \(m\), and all that remains is discounting.

\paragraph{Boundary on the diagonal \(s=m\).}
This is the most important and most characteristic boundary of the
lookback problem.

When \(s=m\), the stock is exactly at its running maximum.
If the stock moves upward, then the running maximum moves upward with
it.
Thus, along the diagonal, a small increase in \(s\) must be accompanied
by the same increase in \(m\).
\paragraph{Boundary condition on the diagonal \(s=m\).}
On the diagonal, the stock is at its running maximum. Over a short time
interval \([t,t+\Delta t]\), if the stock moves upward, then the running
maximum also increases. More precisely, starting from \(S_t=M_t=m\), one
has
\[
M_{t+\Delta t}-m \approx (S_{t+\Delta t}-S_t)^+.
\]
Since
\[
S_{t+\Delta t}-S_t = O(\sqrt{\Delta t}),
\]
the increment of the running maximum has expectation of order
\(\sqrt{\Delta t}\). Writing the dynamic programming relation and
expanding \(u(t+\Delta t,S_{t+\Delta t},M_{t+\Delta t})\) around
\((t,m,m)\), the term involving \(\partial_m u(t,m,m)\) would therefore
produce a contribution of order \(\sqrt{\Delta t}\). Such a term must
vanish in order to obtain a consistent backward PDE. Hence the diagonal
boundary condition is
\[
\boxed{
\partial_m u(t,m,s)\vert_{s = m} = 0,
\qquad t<T.
}
\]

\paragraph{Upper boundary in \(m\).}
In numerical implementations, one usually truncates the \(m\)-domain at
some large value \(M_{\max}\).
Then one must impose an artificial boundary condition there, for
instance by using the asymptotic behavior
\[
u(t,s,m)\approx m-Ke^{-r(T-t)}
\]
for large \(m\), or by choosing a sufficiently large truncation so that
the numerical error is negligible.

\paragraph{Summary of the main conditions.}
The pricing problem is therefore characterized by:
\begin{itemize}[leftmargin=1.4em]
    \item the PDE in the interior \(0<s<m\),
    \item the terminal condition
    \[
    u(T,s,m)=(m-K)^+,
    \]
    \item the boundary condition at \(s=0\),
    \[
    u(t,0,m)=e^{-r(T-t)}(m-K)^+,
    \]
    \item and the diagonal condition at \(s=m\),
    \[
    \partial_m u(t,m,m)=0.
    \]
\end{itemize}

\subsubsection{Interpretation of the additional state variable}

The variable \(m\) plays a very special role in the problem.

\paragraph{A memory variable.}
The running maximum is not an independent market factor like a second
Brownian motion or a stochastic volatility process.
Rather, it is a \emph{memory variable}: it records the relevant part of
the past trajectory of the stock.

This is what allows us to convert a path-dependent problem into a
Markovian one.

\paragraph{Why this variable is enough.}
Although the full stock path \((S_u)_{0\le u\le t}\) contains infinitely
much information, the payoff of the lookback option depends on the past
only through the maximum already achieved.
Therefore the single variable
\[
M_t=\max_{0\le u\le t}S_u
\]
is sufficient to summarize the relevant history.

This is why the augmented state
\[
(t,S_t,M_t)
\]
is enough for pricing.

\paragraph{How \(m\) changes the geometry of the problem.}
For a European option, the state space is simply the half-line
\[
s\in \R_+.
\]
For the lookback option, the state space becomes the wedge
\[
\{(s,m): 0\le s\le m\}.
\]

Thus the introduction of the running maximum enlarges the dimension of
the problem and changes the geometry of the PDE domain.

\paragraph{How \(m\) changes the hedge.}
The additional state variable also affects the hedge.
The price is now
\[
u(t,s,m),
\]
so the delta is
\[
\Delta_t=\partial_s u(t,S_t,M_t).
\]

Thus the hedge ratio depends not only on the current stock level, but
also on the current running maximum.
This already suggests that hedging a lookback option will be more
delicate than hedging a vanilla European call.

\paragraph{Main lesson.}
The additional state variable is the mathematical expression of path
dependence.
It transforms a non-Markovian payoff into a Markovian pricing problem,
but at the cost of increasing the dimension of the state space and
making the pricing and hedging problem more complex.
\subsection{Numerical implementation}

We now discuss how the pricing problem can be implemented numerically.

Compared with the European call case, the main additional difficulty is
that the payoff depends on the running maximum
\[
M_T=\max_{0\le u\le T}S_u.
\]
Therefore, the numerical method must not only simulate or approximate
the stock price, but also keep track of the relevant path functional.

There are two natural numerical viewpoints:
\begin{itemize}[leftmargin=1.4em]
    \item a Monte Carlo viewpoint, where one simulates stock paths and
    tracks the running maximum pathwise;
    \item a PDE viewpoint, where one solves a pricing equation in the
    enlarged state space \((t,s,m)\).
\end{itemize}

For this course, Monte Carlo is often the most intuitive starting point,
because the lookback payoff is directly expressed in terms of the path.
The PDE viewpoint is still important conceptually, because it shows how
path dependence changes the state space and the boundary conditions.

\subsubsection{Choice of numerical method}

For a lookback option, there is no unique numerical method that is
always best.
The choice depends on the objective:
\begin{itemize}[leftmargin=1.4em]
    \item if one wants a simple and flexible pricing tool, Monte Carlo is
    very natural;
    \item if one wants a full pricing surface and numerical Greeks, a PDE
    approach may be preferable;
    \item if one wants to study hedging path by path, Monte Carlo is
    again the natural framework.
\end{itemize}

\paragraph{Why Monte Carlo is natural here.}
The payoff of the fixed-strike lookback call is
\[
(M_T-K)^+.
\]
This payoff is defined directly in terms of the realized stock path.
Therefore, once one can simulate a path
\[
(S_{t_n})_{0\le n\le N},
\]
it is straightforward to compute the discrete running maximum and then
the payoff.

This makes Monte Carlo particularly convenient for:
\begin{itemize}[leftmargin=1.4em]
    \item pricing,
    \item pathwise analysis,
    \item and hedging P\&L comparisons.
\end{itemize}

\paragraph{Why PDEs are still useful.}
The PDE formulation remains important for two reasons.

First, it shows how path dependence can still be treated in a Markovian
framework once the state is enlarged to
\[
(t,S_t,M_t).
\]

Second, the PDE viewpoint is closer to the logic used in the European
option example:
\begin{itemize}[leftmargin=1.4em]
    \item derive the pricing equation,
    \item solve it numerically,
    \item then recover the delta by differentiation.
\end{itemize}

Thus, even if Monte Carlo is often the easiest implementation route for
a lookback payoff, the PDE formulation remains very informative.

\paragraph{Our practical strategy.}
For this example, a natural pedagogical strategy is:
\begin{itemize}[leftmargin=1.4em]
    \item use Monte Carlo as the main pricing tool,
    \item explain clearly how the running maximum is tracked,
    \item and keep the finite-difference viewpoint as a conceptual
    extension of the PDE approach already used for the European option.
\end{itemize}

\subsubsection{Monte Carlo pricing under the pricing model}

Assume first that the stock follows the Black--Scholes dynamics under
the risk-neutral measure:
\[
\mathrm{d}S_t=rS_t\,\mathrm{d}t+\sigma S_t\,\mathrm{d}W_t.
\]

Then the price of the fixed-strike lookback call is
\[
u(0,S_0,M_0)
=
\E^\Qbb\!\left[
e^{-rT}(M_T-K)^+
\right].
\]

\paragraph{Discrete-time approximation.}
To implement this numerically, we choose a time grid
\[
t_n=n\Delta t,
\qquad
\Delta t=\frac{T}{N},
\qquad
n=0,\dots,N.
\]

We simulate the stock path on this grid, for instance with the exact
lognormal scheme
\[
S_{n+1}
=
S_n\exp\!\Big(
\big(r-\tfrac12\sigma^2\big)\Delta t+\sigma\sqrt{\Delta t}\,Z_n
\Big),
\]
where \(Z_n\sim\mathcal N(0,1)\) are independent.

We then define the discrete running maximum by
\[
M_N^{\mathrm{disc}}
:=
\max_{0\le n\le N}S_n.
\]

The Monte Carlo estimator of the option price is
\[
\widehat u_0
=
e^{-rT}\frac1M
\sum_{j=1}^M
\Big(M_{N,j}^{\mathrm{disc}}-K\Big)^+,
\]
where \(M\) is the number of simulated paths.

\paragraph{Interpretation.}
The logic is exactly the same as in standard Monte Carlo pricing:
\begin{itemize}[leftmargin=1.4em]
    \item simulate independent sample paths under the risk-neutral
    measure,
    \item compute the payoff along each path,
    \item average the discounted payoffs.
\end{itemize}

The new feature is only that the payoff is computed from the path
maximum rather than from the terminal stock value alone.

\paragraph{Non-vectorized Python implementation.}
A first direct implementation is as follows.

\begin{lstlisting}[language=Python]
import numpy as np

def simulate_bs_lookback_path_non_vectorized(S0, r, sigma, T, N):
    dt = T / N
    S = np.zeros(N + 1)
    S[0] = S0

    for n in range(N):
        z = np.random.randn()
        S[n + 1] = S[n]*np.exp((r-0.5*sigma**2)*dt+sigma*np.sqrt(dt)*z)

    M = np.max(S)
    return S, M

def price_lookback_mc_non_vectorized(S0, K, r, sigma, T, N, M_paths):
    payoffs = np.zeros(M_paths)

    for j in range(M_paths):
        S_path,M_path=simulate_bs_lookback_path_non_vectorized(S0,r,sigma,T,N)
        payoffs[j] = max(M_path - K, 0.0)

    price = np.exp(-r * T) * np.mean(payoffs)
    return price
\end{lstlisting}

\paragraph{Vectorized Python implementation.}
The same idea can be written much more efficiently in vectorized form.

\begin{lstlisting}[language=Python]
import numpy as np

def simulate_bs_lookback_paths_vectorized(S0, r, sigma, T, N, M_paths):
    dt = T / N
    Z = np.random.randn(M_paths, N)

    log_increments = (r - 0.5 * sigma**2) * dt + sigma * np.sqrt(dt) * Z
    log_paths = np.cumsum(log_increments, axis=1)

    S = np.zeros((M_paths, N + 1))
    S[:, 0] = S0
    S[:, 1:] = S0 * np.exp(log_paths)

    M_running = np.max(S, axis=1)
    return S, M_running

def price_lookback_mc_vectorized(S0, K, r, sigma, T, N, M_paths):
    S,M_running=simulate_bs_lookback_paths_vectorized(S0,r,sigma,T,N,M_paths)
    payoffs = np.maximum(M_running - K, 0.0)
    price = np.exp(-r * T) * np.mean(payoffs)
    return price
\end{lstlisting}

\paragraph{Confidence interval.}
As in any Monte Carlo method, one may also compute an approximate
confidence interval based on the empirical standard deviation of the
discounted payoffs.

If
\[
Y_j=e^{-rT}(M_{N,j}^{\mathrm{disc}}-K)^+,
\]
then an approximate \(95\%\) confidence interval is
\[
\widehat u_0
\pm
1.96\,
\frac{\widehat{\mathrm{std}}(Y)}{\sqrt{M}}.
\]

This is useful when comparing prices or hedging strategies numerically.

\subsubsection{Tracking the running maximum / minimum}

The crucial numerical step for a lookback option is the computation of
the running extremum.

\paragraph{Running maximum on a discrete grid.}
Given a simulated stock path
\[
(S_{t_n})_{0\le n\le N},
\]
the discrete running maximum is defined by
\[
M_n=\max_{0\le k\le n}S_{t_k}.
\]

Thus it can be updated recursively as
\[
M_{n+1}=\max(M_n,S_{t_{n+1}}).
\]

Similarly, if one were working with a running minimum, one would define
\[
m_n=\min_{0\le k\le n}S_{t_k},
\qquad
m_{n+1}=\min(m_n,S_{t_{n+1}}).
\]

\paragraph{Why this is numerically simple but conceptually important.}
At the algorithmic level, updating a running maximum is easy.
However, this simple recursion is the numerical expression of the path
dependence of the contract.

The state variable has changed from:
\[
S_n
\]
to
\[
(S_n,M_n).
\]

This is why path-dependent products are more expensive to price and
hedge numerically than standard European options.

\paragraph{Non-vectorized implementation.}
\begin{lstlisting}[language=Python]
def running_max_non_vectorized(S_path):
    N = len(S_path)
    M_path = np.zeros(N)
    M_path[0] = S_path[0]

    for n in range(1, N):
        M_path[n] = max(M_path[n - 1], S_path[n])

    return M_path
\end{lstlisting}

\paragraph{Vectorized implementation.}
In NumPy, the running maximum over each path can be computed very
efficiently using \texttt{maximum.accumulate}.

\begin{lstlisting}[language=Python]
def running_max_vectorized(S_paths):
    return np.maximum.accumulate(S_paths, axis=1)
\end{lstlisting}

\paragraph{Discretization bias.}
A very important practical issue is that the true running maximum of a
continuous-time process is not exactly equal to the maximum observed on
a discrete simulation grid.

Indeed, the stock may exceed the observed discrete maximum between two
sampling dates and come back down before the next date.
Therefore, the discrete-time approximation
\[
\max_{0\le n\le N}S_{t_n}
\]
typically underestimates the true continuous-time maximum.

As a consequence, the Monte Carlo price of a lookback option computed on
a coarse grid is usually biased downward.

\paragraph{Implication.}
This means that path-dependent products are particularly sensitive to
time discretization.
To reduce the bias, one may:
\begin{itemize}[leftmargin=1.4em]
    \item refine the time grid,
    \item use Brownian-bridge corrections,
    \item or use specialized schemes tailored to extrema.
\end{itemize}

For a first numerical implementation, refining the grid is usually the
simplest approach.

\subsubsection{Finite-difference viewpoint in the extended state space}

Although Monte Carlo is often the most natural implementation tool for a
lookback payoff, it is useful to understand how the PDE point of view
extends to this setting.

\paragraph{The pricing PDE.}
For the fixed-strike lookback call under Black--Scholes, the price
\(u(t,s,m)\) solves in the interior region \(0<s<m\) the PDE
\[
\partial_t u
+
rs\,\partial_s u
+
\frac12 \sigma^2 s^2\,\partial_{ss}u
-r\,u
=
0,
\]
with terminal condition
\[
u(T,s,m)=(m-K)^+,
\]
and suitable boundary conditions at \(s=0\) and \(s=m\).

\paragraph{Main numerical difference with the European case.}
For a European option, the PDE is solved on a two-dimensional grid
\[
(t,s).
\]
For a lookback option, the PDE must be solved on a three-dimensional
grid
\[
(t,s,m),
\qquad 0\le s\le m.
\]

Thus the state space has gained one additional dimension.

\paragraph{Grid structure.}
A finite-difference method would require:
\begin{itemize}[leftmargin=1.4em]
    \item a time grid \(t_n\),
    \item a stock grid \(s_i\),
    \item a maximum grid \(m_j\),
    \item and a price array
    \[
    U[n,i,j]\approx u(t_n,s_i,m_j).
    \]
\end{itemize}

The main complication is that the domain is triangular:
\[
s_i\le m_j.
\]
So not all pairs \((s_i,m_j)\) are admissible.

\paragraph{Interior finite-difference scheme.}
Inside the region \(s<m\), the PDE has no derivative with respect to
\(m\), because the running maximum remains frozen there.
Thus the interior step resembles a Black--Scholes update in the stock
variable, performed for each fixed level of \(m\).

In other words, for each fixed \(m_j\), one solves a one-dimensional PDE
in the stock direction, but coupled through the boundary condition on
the diagonal \(s=m\).

\paragraph{Diagonal boundary.}
The most delicate part of the PDE implementation is the diagonal
boundary \(s=m\), where the stock is at its running maximum.
This is where the geometry of the state space matters.

Numerically, one must impose the compatibility condition derived in the
continuous theory, for instance
\[
\partial_m u(t,m,m)=0,
\]
or an equivalent discrete relation.

This is one reason why PDE methods for lookback options are more
delicate than for vanilla options.

\paragraph{Why PDEs are still worth mentioning.}
Even if one ultimately prices the lookback option by Monte Carlo, the
finite-difference viewpoint remains important because:
\begin{itemize}[leftmargin=1.4em]
    \item it makes explicit the role of the enlarged state space,
    \item it shows how path dependence enters the pricing equation,
    \item and it clarifies how the delta hedge can in principle be
    obtained from a grid-based pricing function.
\end{itemize}

\paragraph{Main practical message.}
For path-dependent products such as lookbacks:
\begin{itemize}[leftmargin=1.4em]
    \item Monte Carlo is often the most direct method for pricing and for
    pathwise hedging studies,
    \item while PDE methods remain very informative conceptually, but
    become more expensive and more delicate because of the enlarged state
    space and the special boundary geometry.
\end{itemize}
\subsection{Hedging the lookback option}

We now turn to the hedging problem.
As in the European-option case, the general idea is to construct a
self-financing portfolio in the stock and the money market account, and
to choose the stock position so as to offset the local exposure of the
option price to stock movements.

However, the lookback option is more delicate to hedge than a European
call, because its price depends not only on the current stock price, but
also on the running maximum of the stock.
Thus the hedge ratio is no longer a function of \((t,S_t)\) only, but of
\((t,S_t,M_t)\).

\subsubsection{Delta hedge under the pricing model}

Let
\[
u(t,s,m)
\]
denote the price of the fixed-strike lookback call at time \(t\) when
\[
S_t=s,
\qquad
M_t=m.
\]

We assume first that this price is computed under the chosen pricing
model, for instance Black--Scholes or local volatility, and that the
function \(u\) is smooth enough.

\paragraph{Self-financing portfolio.}
We consider a hedging portfolio of the form
\[
\Pi_t=\Delta_t S_t+\beta_t B_t,
\]
where:
\begin{itemize}[leftmargin=1.4em]
    \item \(\Delta_t\) is the number of shares held in the stock,
    \item \(\beta_t\) is the amount invested in the money market account,
    \item \(B_t=e^{rt}\).
\end{itemize}

Assuming the strategy is self-financing, the portfolio dynamics satisfy
\[
\mathrm{d}\Pi_t=\Delta_t\,\mathrm{d}S_t+\beta_t\,\mathrm{d}B_t.
\]

Since
\[
\mathrm{d}B_t=rB_t\,\mathrm{d}t,
\]
this can also be written as
\[
\mathrm{d}\Pi_t
=
\Delta_t\,\mathrm{d}S_t+r(\Pi_t-\Delta_t S_t)\,\mathrm{d}t.
\]

\paragraph{Dynamics of the option price in the interior region.}
Inside the region \(S_t<M_t\), the running maximum stays locally
constant.
Therefore, applying It\^o's formula to \(u(t,S_t,M_t)\) gives
\[
\mathrm{d}u(t,S_t,M_t)
=
\Big(
\partial_t u
+
rS_t\,\partial_s u
+
\frac12 \sigma(t,S_t)^2S_t^2\,\partial_{ss}u
\Big)\mathrm{d}t
+
\sigma(t,S_t)S_t\,\partial_s u\,\mathrm{d}W_t,
\]
where \(\sigma(t,S_t)\) denotes the volatility coefficient of the chosen
pricing model:
\begin{itemize}[leftmargin=1.4em]
    \item \(\sigma(t,s)\equiv \sigma\) in Black--Scholes,
    \item \(\sigma(t,s)=\sigma_{\mathrm{loc}}(t,s)\) in the local
    volatility model.
\end{itemize}

\paragraph{Matching the Brownian terms.}
If we want the hedging portfolio to replicate the option price, then in
the interior region we must match the diffusion terms of
\(\Pi_t\) and \(u(t,S_t,M_t)\).
This yields
\[
\Delta_t \sigma(t,S_t)S_t
=
\sigma(t,S_t)S_t\,\partial_s u(t,S_t,M_t),
\]
and therefore
\[
\boxed{
\Delta_t=\partial_s u(t,S_t,M_t).
}
\]

Thus the delta hedge for the lookback option is still the derivative of
the price with respect to the stock, but now the price depends on the
additional state variable \(M_t\).

\paragraph{Interpretation.}
This formula shows that the formal hedging rule looks identical to the
European case:
\[
\Delta_t=\partial_s u.
\]
The crucial difference is that \(u\) is no longer a function of only
\((t,s)\), but of \((t,s,m)\).
Hence the hedge ratio now depends on the current running maximum as
well:
\[
\Delta_t=\partial_s u(t,S_t,M_t).
\]

\paragraph{Black--Scholes benchmark versus richer models.}
If the pricing model is Black--Scholes, then the delta is obtained from
the lookback price function
\[
u^{BS}(t,s,m).
\]
If the pricing model is local volatility, then the hedge ratio becomes
\[
\Delta_t^{LV}=\partial_s u^{LV}(t,S_t,M_t).
\]

So, as in the European-option example, the hedge is model-dependent.
But now the impact of model choice is potentially stronger because the
price depends on path information through \(M_t\).

\subsubsection{Why hedging is more delicate than for a European option}

The lookback option is more difficult to hedge than a standard European
call for several related reasons.

\paragraph{First reason: the payoff depends on the whole path.}
A European call depends only on \(S_T\), so its hedge is driven by the
sensitivity of the price to the current stock level.

A lookback option depends on
\[
M_T=\max_{0\le u\le T}S_u,
\]
so the future payoff is affected not only by where the stock ends, but
also by whether and when the stock reaches new maxima before maturity.

Thus the option is sensitive to the full trajectory of the stock.

\paragraph{Second reason: the state variable is richer.}
For a European call, the price is
\[
u(t,s),
\]
and the delta is a function of \((t,s)\).
For a lookback option, the price is
\[
u(t,s,m),
\]
and the delta depends on both \(s\) and \(m\).

So two states with the same current stock price \(s\) but different
running maxima \(m\) will generally lead to different hedge ratios.

\paragraph{Third reason: the behavior near the diagonal is delicate.}
When the stock is close to the running maximum, that is, when
\[
S_t\approx M_t,
\]
a small upward move may generate a new maximum and thereby change the
future payoff structure.

This creates a region where the price function is especially sensitive
to the interaction between the stock and the path variable.
As a consequence, the delta may vary rapidly near the diagonal
\[
s=m.
\]

\paragraph{Fourth reason: discrete hedging is more fragile.}
In continuous time, the model updates the running maximum continuously.
In practice, however, rebalancing is done on a discrete grid of times.
Between two rebalancing dates, the stock may make a brief excursion
above the previous maximum and then come back.
Such a move may have a significant effect on the final payoff, even if
it is not perfectly captured by the discrete hedge.

Thus discretization error is often more severe for path-dependent
products than for standard vanillas.

\paragraph{Fifth reason: model misspecification matters more.}
For a European call, the hedge depends mainly on the marginal behavior
of the stock under the chosen model.
For a lookback option, the hedge depends on the path behavior of the
stock, and therefore on how the model describes:
\begin{itemize}[leftmargin=1.4em]
    \item the frequency of upward excursions,
    \item the probability of reaching new maxima,
    \item the local behavior of volatility along the path.
\end{itemize}

So the hedging performance is typically more sensitive to model choice.

\paragraph{Main message.}
The lookback option is harder to hedge not because the formal hedging
rule changes, but because the underlying price function is more complex,
more path-sensitive, and more model-dependent than in the European case.

\subsubsection{Discrete-time rebalancing}

As in the European-option example, the ideal hedge in continuous time is
only a theoretical benchmark.
In practice, the hedging portfolio is rebalanced at discrete dates
\[
0=t_0<t_1<\cdots<t_N=T.
\]

\paragraph{Discrete hedge ratio.}
At each rebalancing date \(t_n\), we compute the hedge ratio from the
pricing model:
\[
\Delta_{t_n}
=
\partial_s u(t_n,S_{t_n},M_{t_n}).
\]

Thus the strategy must keep track of both:
\begin{itemize}[leftmargin=1.4em]
    \item the current stock price \(S_{t_n}\),
    \item the current running maximum \(M_{t_n}\).
\end{itemize}

\paragraph{Portfolio decomposition.}
Suppose that we sell one lookback option at time \(0\), and use the
initial model price \(u(0,S_0,M_0)\) to set up the hedging portfolio.
The portfolio value is written as
\[
\Pi_{t_n}
=
\Delta_{t_n}S_{t_n}
+
B_{t_n}^{\mathrm{cash}},
\]
where \(B_{t_n}^{\mathrm{cash}}\) is the cash position.

Initially, one may set
\[
\Pi_0=u(0,S_0,M_0),
\qquad
B_0^{\mathrm{cash}}
=
u(0,S_0,M_0)-\Delta_0 S_0.
\]

\paragraph{Update between two dates.}
Between \(t_n\) and \(t_{n+1}\), the stock holding is kept constant, and
the cash position accrues interest.
Therefore, just before rebalancing at time \(t_{n+1}\), the portfolio
value is
\[
\Pi_{t_{n+1}}^{-}
=
e^{r\Delta t}
\big(\Pi_{t_n}-\Delta_{t_n}S_{t_n}\big)
+
\Delta_{t_n}S_{t_{n+1}},
\]
where
\[
\Delta t=t_{n+1}-t_n.
\]

\paragraph{Rebalancing at the new date.}
At time \(t_{n+1}\), we update the running maximum:
\[
M_{t_{n+1}}=\max(M_{t_n},S_{t_{n+1}}),
\]
compute the new hedge ratio
\[
\Delta_{t_{n+1}}=\partial_s u(t_{n+1},S_{t_{n+1}},M_{t_{n+1}}),
\]
and redefine the cash position so that
\[
\Pi_{t_{n+1}}
=
\Delta_{t_{n+1}}S_{t_{n+1}}
+
B_{t_{n+1}}^{\mathrm{cash}}.
\]

\paragraph{What is new compared with the European case?}
The portfolio update itself is formally the same as before.
The crucial difference is that the hedge ratio is no longer determined
by \((t_n,S_{t_n})\) alone, but by the augmented state
\[
(t_n,S_{t_n},M_{t_n}).
\]

So, at each date, the hedger must update not only the stock position and
the cash account, but also the running maximum used in the pricing and
hedging formula.

\paragraph{Practical implication.}
This makes the hedge more sensitive to path recording and to the
accuracy of the numerical scheme.
Any error in tracking the running maximum may translate into an error in
the hedge ratio.

\subsubsection{Sensitivity to the path and to the proximity of the maximum / minimum}

The lookback option is especially sensitive to where the current stock
stands relative to the running maximum.

\paragraph{When the stock is far below the maximum.}
If
\[
S_t\ll M_t,
\]
then the option payoff is already strongly influenced by the previously
reached maximum.
In that situation, the probability that a new higher maximum will be
created before maturity may be relatively small.
The hedge is then less driven by immediate stock moves and more by the
fact that a large historical maximum is already locked in.

\paragraph{When the stock is close to the maximum.}
If
\[
S_t\approx M_t,
\]
then even a small upward move may create a new maximum.
This can materially affect the future payoff
\[
(M_T-K)^+.
\]

As a result, the option price may become very sensitive to local stock
movements in this region, and the delta may change rapidly.
This is one of the reasons why hedging is particularly delicate near the
diagonal
\[
s=m.
\]

\paragraph{Interpretation in terms of path sensitivity.}
The current value of the running maximum summarizes the relevant history
of the stock path.
The closer the stock is to this historical maximum, the more the future
payoff is sensitive to the next local move of the stock.

So the hedge is not determined only by the current level of the stock,
but by the interaction between:
\begin{itemize}[leftmargin=1.4em]
    \item the current stock value,
    \item the historical maximum already achieved,
    \item and the remaining time to maturity.
\end{itemize}

\paragraph{Extension to the running minimum.}
If the contract depended on the running minimum instead of the running
maximum, the same logic would apply in reverse.
The hedge would then be particularly sensitive when the stock is close
to the running minimum.

Thus, for path-dependent products based on extrema, proximity to the
relevant boundary is a key driver of hedging sensitivity.

\paragraph{Why this matters for model risk.}
Any pricing model determines how likely the stock is to create new
maxima or minima.
Since the lookback hedge is highly sensitive to those events, a wrong
model may produce a significantly wrong hedge ratio, especially near the
extremal boundary.

This makes the lookback option an excellent example for studying the
interaction between:
\begin{itemize}[leftmargin=1.4em]
    \item path dependence,
    \item numerical approximation,
    \item and model misspecification.
\end{itemize}

\paragraph{Main lesson.}
The hedge of a lookback option is more fragile than the hedge of a
vanilla option because it depends on the whole relevant history of the
path through the running extremum.
This path sensitivity becomes especially strong when the stock is close
to the current maximum or minimum.

\subsection{Model risk for path-dependent products}

We now return to the issue of model risk, but in the more demanding
setting of a path-dependent product.

For the European call, model risk was already visible: two different
hedging models could lead to different deltas and different hedging
P\&L distributions when the true dynamics were stochastic volatility.

For a lookback option, this effect is typically even stronger.
The reason is that the payoff depends on the whole path of the stock,
and therefore on the way the model describes the pathwise evolution of
volatility, the probability of reaching new extrema, and the frequency
of large excursions.

\subsubsection{A reference stochastic volatility model for the true dynamics}

To study model risk concretely, we assume that the \emph{true} stock
dynamics are given by a stochastic volatility model, while the hedging
strategies are built under simpler pricing models.

A natural reference choice is a Heston-type model under the
risk-neutral measure:
\[
\mathrm{d}S_t
=
rS_t\,\mathrm{d}t
+
\sqrt{V_t}\,S_t\,\mathrm{d}W_t^{(1)},
\]
\[
\mathrm{d}V_t
=
\kappa(\theta-V_t)\,\mathrm{d}t
+
\eta\sqrt{V_t}\,\mathrm{d}W_t^{(2)},
\]
with
\[
\mathrm{d}\langle W^{(1)},W^{(2)}\rangle_t=\rho\,\mathrm{d}t.
\]

Here:
\begin{itemize}[leftmargin=1.4em]
    \item \(S_t\) is the stock price,
    \item \(V_t\) is the instantaneous variance,
    \item \(\kappa\) is the mean-reversion speed,
    \item \(\theta\) is the long-run variance level,
    \item \(\eta\) is the volatility of volatility,
    \item \(\rho\) is the correlation between stock and variance shocks.
\end{itemize}

\paragraph{Why this is a natural benchmark.}
This model is rich enough to produce:
\begin{itemize}[leftmargin=1.4em]
    \item time-varying volatility,
    \item volatility clustering,
    \item asymmetric smile effects when \(\rho\neq 0\),
    \item and path behavior that differs substantially from a constant-
    volatility diffusion.
\end{itemize}

Thus it provides a meaningful ``true world'' against which simpler
hedging models can be tested.

\paragraph{Risk-neutral pricing under the true model.}
If we were to price the lookback option directly under this stochastic
volatility model, the natural risk-neutral value would be
\[
u^{SV}(t,s,v,m)
=
\E^\Qbb\!\left[
e^{-r(T-t)}(M_T-K)^+
\;\middle|\;
S_t=s,\ V_t=v,\ M_t=m
\right].
\]

So the true Markovian state variable would now be
\[
(t,S_t,V_t,M_t),
\]
which is even richer than the state space used under the benchmark
Black--Scholes or local-volatility pricing models.

\paragraph{Why we do not necessarily hedge under the true model.}
In practice, a trader rarely knows the true model.
The whole point of model risk is precisely that the true dynamics are
unknown.
So in our numerical experiment, the stochastic volatility model is used
only to generate the true stock paths and the true running maximums,
while the hedging strategies themselves are computed under simpler
models.

This is exactly the right setup for illustrating model misspecification.

\subsubsection{Why model misspecification matters more here}

Model misspecification already affects the hedge of a European option,
but it typically matters even more for a path-dependent product such as
a lookback option.

\paragraph{European case versus lookback case.}
For a European call, the payoff depends only on the terminal stock
value:
\[
(S_T-K)^+.
\]
So pricing and hedging depend essentially on the model's description of
the terminal distribution of \(S_T\).

For a lookback option, the payoff depends on the running maximum:
\[
(M_T-K)^+,
\qquad
M_T=\max_{0\le u\le T}S_u.
\]
Thus pricing and hedging depend not only on where the stock ends, but
also on how it moves during the whole interval \([0,T]\).

\paragraph{Why the path matters.}
Two models may produce similar prices for liquid European options at
time \(0\), yet generate very different path behavior for the stock.
In particular, they may differ in:
\begin{itemize}[leftmargin=1.4em]
    \item how often the stock approaches its running maximum,
    \item how likely it is to break through an existing maximum,
    \item how volatility behaves during large upward or downward moves,
    \item how extreme path realizations are distributed over time.
\end{itemize}

These differences can have a strong effect on a lookback payoff.

\paragraph{Why stochastic volatility is particularly relevant.}
In a stochastic volatility model, the variance itself evolves randomly.
This means that some paths may experience:
\begin{itemize}[leftmargin=1.4em]
    \item periods of low volatility,
    \item periods of high volatility,
    \item bursts of volatility coinciding with large stock moves.
\end{itemize}

Such effects may significantly alter the probability that the stock
reaches new maxima before maturity.

A simple Black--Scholes or local-volatility hedging model does not
reproduce these dynamics in the same way.
Therefore the resulting hedge may be significantly misspecified.

\paragraph{Path dependence amplifies model risk.}
This is the key intuition:
\begin{quote}
for a path-dependent product, a model error affects not only the final
distribution of the stock, but the whole evolution of the path, and
therefore the whole mechanism that generates the payoff.
\end{quote}

As a result, the impact of model misspecification is often amplified for
lookback options compared with standard European calls.

\paragraph{Another source of fragility: the hedge depends on the running maximum.}
The hedge ratio of a lookback option depends on
\[
(t,S_t,M_t).
\]
So any model error that changes the distribution of future maxima also
changes the relevance of the hedge ratio computed under the wrong model.

This is especially important near the region where
\[
S_t\approx M_t,
\]
because in that region the hedge is very sensitive to the probability of
creating a new maximum.

\paragraph{Main message.}
For a path-dependent product, model misspecification matters more
because the model must get the \emph{path dynamics} reasonably right,
not only the terminal law.
This is why model risk is particularly visible in the hedging P\&L of
lookback options.

\subsubsection{Hedging with a simple model versus hedging with a richer model}

We now compare two ways of building a hedge for the lookback option when
the true dynamics are stochastic volatility.

\paragraph{A simple hedging model.}
A first approach is to hedge under a simple benchmark model, for
instance Black--Scholes:
\[
\mathrm{d}S_t=rS_t\,\mathrm{d}t+\sigma S_t\,\mathrm{d}W_t.
\]

In this case, the hedge ratio is computed from the benchmark lookback
price
\[
u^{BS}(t,s,m),
\]
and the corresponding delta is
\[
\Delta_t^{BS}
=
\partial_s u^{BS}(t,S_t,M_t).
\]

This strategy is easy to interpret and to implement.
It provides a clean benchmark.
But it ignores:
\begin{itemize}[leftmargin=1.4em]
    \item smile effects,
    \item state dependence of volatility,
    \item and random fluctuations of the true variance process.
\end{itemize}

\paragraph{A richer hedging model.}
A second approach is to hedge under a richer model, such as local
volatility:
\[
\mathrm{d}S_t
=
rS_t\,\mathrm{d}t
+
\sigma_{\mathrm{loc}}(t,S_t)S_t\,\mathrm{d}W_t.
\]

The hedge ratio is then computed from the local-volatility lookback
price
\[
u^{LV}(t,s,m),
\]
and takes the form
\[
\Delta_t^{LV}
=
\partial_s u^{LV}(t,S_t,M_t).
\]

This strategy is still simpler than a full stochastic-volatility hedge,
but it is at least consistent with the current market smile through the
local volatility surface.

\paragraph{Why the richer model may help.}
Even if local volatility is not the true model, it may provide a better
hedge than Black--Scholes because it uses more information from liquid
market vanillas.
In particular, it may produce stock-path dynamics that are more
consistent with observed smile effects than the constant-volatility
benchmark.

Thus it may lead to a hedge ratio that is more robust when the true
world is not Black--Scholes.

\paragraph{Why the richer model may still fail.}
However, local volatility is still not the true stochastic volatility
model.
It incorporates the current smile, but volatility remains a deterministic
function of \((t,S_t)\).
It does not have a separate random volatility factor.

So even the richer hedge remains misspecified when the true dynamics are
stochastic volatility.

\paragraph{What we want to compare numerically.}
The numerical experiment therefore compares:
\begin{itemize}[leftmargin=1.4em]
    \item a simple hedge built under Black--Scholes,
    \item a richer hedge built under local volatility,
    \item both tested on paths generated under stochastic volatility.
\end{itemize}

This allows us to ask a concrete and practically meaningful question:
\begin{quote}
does a richer hedging model produce a smaller and more stable hedging
error when the true dynamics are pathwise more complex?
\end{quote}

\paragraph{What we should not expect.}
We should not expect either hedge to be perfect.
Since the true dynamics are stochastic volatility, both the
Black--Scholes hedge and the local-volatility hedge are misspecified.
So the purpose of the comparison is not to find an exact replication
strategy, but to compare the size and distribution of the residual
hedging error.

\paragraph{Main lesson.}
For a path-dependent product such as a lookback option, one expects the
difference between a simple hedge and a richer hedge to be more visible
than for a vanilla European option.
The reason is that the product is more sensitive to the detailed
pathwise behavior of the stock, and therefore more sensitive to model
choice.

\subsection{Comparing hedging P\&L}

We now explain how to compare the performance of the hedging strategies
for the lookback option.

As in the European-option example, the idea is to:
\begin{itemize}[leftmargin=1.4em]
    \item choose a pricing and hedging model,
    \item simulate the \emph{true} stock dynamics under a richer model,
    \item implement the hedge in discrete time,
    \item and study the terminal profit and loss.
\end{itemize}

For a path-dependent product, this comparison is especially informative,
because the hedge depends not only on the stock level, but also on the
running maximum, and therefore on the whole realized path.

\subsubsection{Definition of the hedging error}

Suppose that we sell one fixed-strike lookback call at time \(0\), with
payoff
\[
(M_T-K)^+,
\qquad
M_T=\max_{0\le u\le T}S_u.
\]

Assume that the option is priced under the chosen hedging model, giving
the initial price
\[
u(0,S_0,M_0).
\]

We use this initial amount to set up a self-financing hedging
portfolio, which is then rebalanced at discrete dates
\[
0=t_0<t_1<\cdots<t_N=T.
\]

Let \(\Pi_T\) denote the final value of this hedging portfolio at
maturity.

\paragraph{Terminal hedging P\&L.}
Since we are short one option, the terminal hedging profit and loss is
defined by
\[
\boxed{
\mathrm{PnL}
=
\Pi_T-(M_T-K)^+.
}
\]

\paragraph{Interpretation.}
If the hedge were perfect and rebalancing continuous in the true model,
then one would have
\[
\Pi_T=(M_T-K)^+,
\]
so that
\[
\mathrm{PnL}=0.
\]

In practice, however, the terminal P\&L is generally not zero because of
several sources of error:
\begin{itemize}[leftmargin=1.4em]
    \item discrete rebalancing,
    \item numerical approximation,
    \item and, most importantly here, model misspecification.
\end{itemize}

\paragraph{Why the notion of error is more subtle here.}
For a European option, the payoff depends only on \(S_T\), so the
hedging error is driven mainly by how well the model captures the
terminal exposure of the option.

For a lookback option, the payoff depends on the running maximum over
the whole interval \([0,T]\).
Thus the hedging error reflects not only the final stock level, but also
how well the model captures the full path behavior of the stock.

\paragraph{Two candidate hedging errors.}
In our numerical comparison, we typically define:
\begin{itemize}[leftmargin=1.4em]
    \item one P\&L associated with the Black--Scholes hedge,
    \item one P\&L associated with the local-volatility hedge.
\end{itemize}

If \(\Pi_T^{BS}\) and \(\Pi_T^{LV}\) denote the final values of the two
hedging portfolios, then
\[
\mathrm{PnL}^{BS}
=
\Pi_T^{BS}-(M_T-K)^+,
\]
\[
\mathrm{PnL}^{LV}
=
\Pi_T^{LV}-(M_T-K)^+.
\]

The comparison of these two random variables is the central object of
this part of the analysis.

\subsubsection{Pathwise comparison of hedging strategies}

A first level of analysis is to compare the two hedging strategies
\emph{path by path}.

\paragraph{What is done pathwise.}
For each simulated trajectory of the true stock dynamics, we can:
\begin{itemize}[leftmargin=1.4em]
    \item compute the realized stock path \((S_{t_n})_{0\le n\le N}\),
    \item compute the corresponding running maximum path
    \((M_{t_n})_{0\le n\le N}\),
    \item construct the Black--Scholes hedge along that path,
    \item construct the local-volatility hedge along that path,
    \item compute the final payoff \((M_T-K)^+\),
    \item and evaluate the terminal P\&L for each strategy.
\end{itemize}

\paragraph{Why pathwise comparison is useful.}
Pathwise comparison is particularly important for a lookback option,
because the difference between the two strategies may be driven by very
specific events along the path, for example:
\begin{itemize}[leftmargin=1.4em]
    \item a sudden approach to the running maximum,
    \item the creation of a new historical maximum,
    \item a period of elevated volatility,
    \item or a sequence of oscillations near the maximum.
\end{itemize}

Such events may have a moderate impact on a vanilla European option, but
a much stronger effect on a path-dependent product.

\paragraph{A typical pathwise scenario.}
Suppose that two hedging models produce similar deltas when the stock is
far below the running maximum.
Then along such parts of the trajectory, the two portfolios may evolve
similarly.

However, if at some date the stock comes close to the running maximum,
then the local-volatility hedge and the Black--Scholes hedge may react
very differently.
This may lead to significantly different portfolio values by maturity,
even though both strategies started from the same initial contract.

\paragraph{What one may plot.}
A pathwise comparison is often best illustrated by plotting, for one or
several simulated paths:
\begin{itemize}[leftmargin=1.4em]
    \item the stock path \(S_t\),
    \item the running maximum \(M_t\),
    \item the two hedge ratios \(\Delta_t^{BS}\) and \(\Delta_t^{LV}\),
    \item the cumulative values of the two hedging portfolios,
    \item and the final realized payoff.
\end{itemize}

These plots help visualize where the two hedges begin to diverge and why
their final P\&L may differ.

\paragraph{Main pathwise intuition.}
For a path-dependent product, the hedging error is often produced by a
small number of critical path events rather than by a smooth accumulation
of small local errors.
This is one reason why the analysis of hedging P\&L is especially
interesting for lookback options.

\subsubsection{Distribution of the hedging P\&L}

Pathwise analysis is useful qualitatively, but to compare the two
strategies in a systematic way we need to study the full distribution of
their terminal P\&L.

\paragraph{Monte Carlo sample of hedging errors.}
After simulating many independent paths of the true stochastic
volatility model, we obtain:
\begin{itemize}[leftmargin=1.4em]
    \item a sample of values of \(\mathrm{PnL}^{BS}\),
    \item a sample of values of \(\mathrm{PnL}^{LV}\).
\end{itemize}

These samples can be analyzed statistically.

\paragraph{Natural summary statistics.}
Useful quantities to compare include:
\begin{itemize}[leftmargin=1.4em]
    \item the empirical mean,
    \item the empirical standard deviation,
    \item lower and upper quantiles,
    \item the probability of large negative P\&L,
    \item and more generally the overall concentration around zero.
\end{itemize}

\paragraph{Interpretation of the mean.}
If the empirical mean of the hedging P\&L is significantly different
from zero, this may indicate a systematic bias in the hedge.
Such a bias may come from:
\begin{itemize}[leftmargin=1.4em]
    \item model misspecification,
    \item discretization effects,
    \item or a persistent mismatch between the hedging model and the
    true path dynamics.
\end{itemize}

\paragraph{Interpretation of the variance.}
The empirical variance measures how volatile the residual hedging error
is.
A smaller variance means that the hedge is more stable across scenarios,
which is generally desirable.

\paragraph{Why quantiles matter.}
For risk management, the tails of the P\&L distribution are often more
important than the mean.
A hedge that performs well on average but occasionally produces very
large losses may be less attractive than a hedge with a slightly larger
average error but much smaller tail risk.

This is especially relevant for path-dependent products, where rare path
events may generate large hedging losses.

\paragraph{Visual comparison.}
The distribution of the hedging P\&L can be displayed using:
\begin{itemize}[leftmargin=1.4em]
    \item histograms,
    \item kernel density estimates,
    \item boxplots,
    \item or empirical cumulative distribution functions.
\end{itemize}

These visual tools make it easy to compare:
\begin{itemize}[leftmargin=1.4em]
    \item the center of the distribution,
    \item the spread,
    \item and the shape of the tails.
\end{itemize}

\paragraph{What we are really comparing.}
We are not only comparing two numbers.
We are comparing two full distributions:
\[
\mathrm{Law}(\mathrm{PnL}^{BS})
\qquad\text{and}\qquad
\mathrm{Law}(\mathrm{PnL}^{LV}).
\]

This is important because one strategy may have:
\begin{itemize}[leftmargin=1.4em]
    \item a smaller mean error,
    \item but a larger dispersion,
\end{itemize}
while another may:
\begin{itemize}[leftmargin=1.4em]
    \item have a slightly larger average error,
    \item but much smaller downside tail risk.
\end{itemize}

Thus the comparison of hedging models is naturally a comparison of risk
profiles, not just of average performance.

\subsubsection{Interpretation of the numerical results}

The numerical results should be interpreted with care.

\paragraph{We do not expect a perfect hedge.}
Since the true dynamics are stochastic volatility, while the hedging
strategies are computed under simpler models such as Black--Scholes or
local volatility, neither hedge is exact.

Thus the purpose of the experiment is not to identify a perfect
replication strategy, but rather to compare the size and nature of the
residual hedging errors.

\paragraph{What one may expect qualitatively.}
Because local volatility incorporates information from the current
vanilla smile, one may expect it to perform better than Black--Scholes
in some scenarios.
This is especially plausible when the lookback option is sensitive to
regions of the state space where the smile matters strongly.

However, this is not guaranteed.
Since the true world is stochastic volatility rather than local
volatility, the local-volatility hedge is itself misspecified.
Its advantage is therefore only partial.

\paragraph{Why the difference may be stronger than for a vanilla option.}
For a European option, the hedge depends mainly on the terminal
distribution of the stock.
For a lookback option, the hedge depends on the whole path, and therefore
on how well the model captures:
\begin{itemize}[leftmargin=1.4em]
    \item the frequency of excursions to new maxima,
    \item the volatility of the path near the running maximum,
    \item and the interaction between path behavior and local hedging.
\end{itemize}

For this reason, the difference between two hedging models is often more
visible for a lookback option than for a vanilla call.

\paragraph{Why no universal ranking should be expected.}
The comparison between Black--Scholes and local volatility depends on:
\begin{itemize}[leftmargin=1.4em]
    \item the strike,
    \item the maturity,
    \item the true stochastic-volatility parameters,
    \item the rebalancing frequency,
    \item and the numerical approximation scheme.
\end{itemize}

Therefore, one should not expect a universal theorem of the form
\[
\text{local volatility hedge is always better than Black--Scholes}.
\]

Rather, the correct interpretation is that richer models may reduce some
forms of model risk, but they do not eliminate model risk entirely.

\paragraph{Main interpretation.}
The key lesson of the numerical experiment is the following:
\begin{quote}
for path-dependent products, model choice matters not only for the
initial price, but even more for the distribution of the future hedging
error.
\end{quote}

This is the practical meaning of model risk in a hedging context.

\subsection{Main lessons from the lookback option example}

The lookback option example highlights several important ideas that go
beyond the standard European-option setting.

\paragraph{First lesson: path dependence changes the state space.}
For a European option, the current stock price is enough to describe the
pricing problem.
For a lookback option, the current stock price is not sufficient.
One must also keep track of the running maximum or running minimum.

Thus the state space becomes larger:
\[
(t,S_t)
\quad\longrightarrow\quad
(t,S_t,M_t).
\]

This is the mathematical expression of path dependence.

\paragraph{Second lesson: path dependence makes pricing and hedging more complex.}
The formal pricing logic remains the same:
\begin{itemize}[leftmargin=1.4em]
    \item choose a model,
    \item write the risk-neutral pricing formula,
    \item derive the PDE,
    \item compute the delta,
    \item and construct a self-financing hedge.
\end{itemize}

But each of these steps becomes more delicate because the payoff depends
on the whole stock path through the running maximum.

\paragraph{Third lesson: the hedge depends on the path history.}
The hedge ratio of a lookback option is no longer just a function of
time and spot.
It depends also on the current running maximum:
\[
\Delta_t=\partial_s u(t,S_t,M_t).
\]

Thus two states with the same stock price but different historical
maxima may require different hedging positions.

\paragraph{Fourth lesson: model risk is amplified for path-dependent products.}
Because the payoff depends on the full path, any model error in the
description of stock-path behavior can have a strong effect on:
\begin{itemize}[leftmargin=1.4em]
    \item the price,
    \item the hedge ratio,
    \item and the final hedging P\&L.
\end{itemize}

This is why model misspecification is often more serious for path-
dependent products than for standard European options.

\paragraph{Fifth lesson: richer models may help, but do not solve everything.}
A richer model such as local volatility may provide a more
market-consistent hedge than the Black--Scholes benchmark, because it
incorporates smile information from liquid vanilla options.

However, if the true dynamics are stochastic volatility, even the local-
volatility hedge remains misspecified.
So the comparison is not between a wrong model and a perfect model, but
between different imperfect hedging models.

\paragraph{Final message.}
Taken together, the European-call example and the lookback-option
example illustrate a central theme of the course:
\[
\text{contract structure}
\longrightarrow
\text{model choice}
\longrightarrow
\text{pricing formula}
\longrightarrow
\text{hedging strategy}
\longrightarrow
\text{hedging P\&L}.
\]

The more path-dependent and model-sensitive the contract is, the more
important this chain becomes.

\paragraph{Transition.}
This concludes the second case study.
The contrast with the European option is now clear:
\begin{itemize}[leftmargin=1.4em]
    \item for a vanilla option, a simple benchmark model already gives a
    very clean pricing and hedging framework;
    \item for a path-dependent option, the state space is richer, the
    hedge is more fragile, and model risk becomes more visible.
\end{itemize}


\section{Option pricing with stochastic interest rates}

\subsection{Why stochastic rates change the problem}

In the standard Black--Scholes framework, the interest rate is usually
assumed to be constant.
This simplification is mathematically very convenient, because it means
that discounting is deterministic:
\[
e^{-r(T-t)}.
\]
As a consequence, the time value of money does not introduce any
additional source of uncertainty.

When interest rates become stochastic, this changes the nature of the
pricing and hedging problem in several important ways.

\paragraph{First change: discounting becomes random.}
If the short rate is now a stochastic process \((r_t)_{t\ge0}\), then
the discount factor between \(t\) and \(T\) becomes
\[
\exp\!\Big(-\int_t^T r_u\,\mathrm{d}u\Big),
\]
which is no longer deterministic.

Therefore, even for a payoff that depends only on the stock price, such
as
\[
(S_T-K)^+,
\]
the discounting part of the pricing formula becomes random.

This is already a major conceptual difference with the constant-rate
case.

\paragraph{Second change: the stock dynamics depend on the short rate.}
Under the risk-neutral measure, the drift of the stock is the short
rate.
So if the short rate is random, then the stock dynamics also change.
Instead of
\[
\mathrm{d}S_t=rS_t\,\mathrm{d}t+\sigma S_t\,\mathrm{d}W_t,
\]
we now have
\[
\mathrm{d}S_t=r_tS_t\,\mathrm{d}t+\cdots.
\]

Thus the random evolution of rates affects not only discounting, but
also the drift of the stock under the pricing measure.

\paragraph{Third change: the state space becomes larger.}
In the constant-rate Black--Scholes model, the option price can be
written as a function
\[
u(t,s),
\]
because the relevant state variables are just time and the current stock
price.

With stochastic rates, the current short rate is itself a state
variable.
So the pricing function naturally becomes
\[
u(t,s,r).
\]

This means:
\begin{itemize}[leftmargin=1.4em]
    \item the pricing problem becomes higher-dimensional,
    \item the PDE becomes more complex,
    \item and the hedge is no longer driven only by the stock
    sensitivity.
\end{itemize}

\paragraph{Fourth change: stock delta alone is no longer enough.}
In the constant-rate case, hedging a European option can often be
explained using only:
\begin{itemize}[leftmargin=1.4em]
    \item the stock,
    \item and the money market account.
\end{itemize}

When rates are stochastic, the option price depends not only on the
stock level, but also on the current interest-rate environment.
So the option is exposed to rate risk as well.

This suggests that, in general, a hedge based only on the stock will not
be sufficient.
One may also need to trade in interest-rate-sensitive instruments such
as zero-coupon bonds or swaps.

\paragraph{Fifth change: the issue of completeness becomes more subtle.}
With stochastic rates, the market may involve several sources of
uncertainty:
\begin{itemize}[leftmargin=1.4em]
    \item one source driving the stock,
    \item one source driving the short rate,
    \item and possibly a correlation between them.
\end{itemize}

Whether or not the option payoff can be perfectly replicated depends on
whether the available traded assets span all these sources of risk.

Thus, compared with the Black--Scholes benchmark, the relation between:
\[
\text{pricing} \quad\text{and}\quad \text{replication}
\]
becomes richer and more delicate.

\paragraph{Main message.}
The introduction of stochastic rates changes the problem because:
\begin{itemize}[leftmargin=1.4em]
    \item discounting is no longer deterministic,
    \item the stock drift under the risk-neutral measure becomes random,
    \item the state space becomes larger,
    \item hedging requires rate-sensitive instruments,
    \item and exact replication is no longer automatic.
\end{itemize}

So, even for a standard option payoff, stochastic rates create a more
realistic but also more complex pricing and hedging problem.

\subsection{A simple market model with stock and short rate}

We now introduce a simple continuous-time market model in which both the
stock price and the short rate evolve randomly.

The purpose of this model is not to capture every feature of fixed-
income markets, but rather to provide the simplest framework in which we
can discuss:
\begin{itemize}[leftmargin=1.4em]
    \item option pricing with stochastic discounting,
    \item the associated PDE,
    \item and the question of hedging and replication.
\end{itemize}

\subsubsection{Risk-neutral dynamics}

We consider a filtered probability space carrying two Brownian motions.
Under the risk-neutral measure \(\Qbb\), we assume that the stock price
\(S_t\) and the short rate \(r_t\) satisfy
\[
\mathrm{d}S_t
=
r_tS_t\,\mathrm{d}t
+
\sigma_S(t,S_t,r_t)\,S_t\,\mathrm{d}W_t^{(1)},
\]
\[
\mathrm{d}r_t
=
b(t,r_t)\,\mathrm{d}t
+
\sigma_r(t,r_t)\,\mathrm{d}W_t^{(2)},
\]
with
\[
\mathrm{d}\langle W^{(1)},W^{(2)}\rangle_t=\rho\,\mathrm{d}t,
\qquad -1\le \rho\le 1.
\]

\paragraph{Interpretation of the stock equation.}
The stock still has drift equal to the short rate under the
risk-neutral measure, exactly as in the Black--Scholes model.
The difference is that the rate is now random.

The volatility coefficient of the stock may depend on:
\begin{itemize}[leftmargin=1.4em]
    \item time,
    \item the current stock level,
    \item and the current short rate.
\end{itemize}

For a first reading, one may keep in mind the simpler specification
\[
\mathrm{d}S_t
=
r_tS_t\,\mathrm{d}t
+
\sigma S_t\,\mathrm{d}W_t^{(1)},
\]
with constant stock volatility \(\sigma\).

\paragraph{Interpretation of the short-rate equation.}
The short rate \(r_t\) is itself modeled as a diffusion.
Its drift is given by a function
\[
b(t,r),
\]
and its diffusion coefficient by
\[
\sigma_r(t,r).
\]

This is the main new ingredient compared with the standard
Black--Scholes framework.

\paragraph{Role of the correlation.}
The correlation parameter \(\rho\) captures the dependence between:
\begin{itemize}[leftmargin=1.4em]
    \item shocks on the stock,
    \item and shocks on the short rate.
\end{itemize}

When \(\rho\neq 0\), stock risk and rate risk are not independent.
This matters for both pricing and hedging, because it creates a mixed
second-order term in the pricing PDE.

\paragraph{Markovian state variable.}
In this model, the natural state variable is
\[
(t,S_t,r_t).
\]
Therefore, the price of a European claim will generally be written as a
function
\[
u(t,s,r).
\]

This is the simplest extension of the Black--Scholes Markovian setting
to the case of stochastic rates.

\paragraph{A useful benchmark simplification.}
For many pedagogical derivations, it is helpful to start with the
simpler model
\[
\mathrm{d}S_t
=
r_tS_t\,\mathrm{d}t
+
\sigma S_t\,\mathrm{d}W_t^{(1)},
\]
\[
\mathrm{d}r_t
=
b(t,r_t)\,\mathrm{d}t
+
\eta(t,r_t)\,\mathrm{d}W_t^{(2)},
\]
where \(\sigma\) is constant.
This keeps the main conceptual novelty, namely stochastic rates, while
avoiding too much notational complexity.

\subsubsection{Money market account}

As in the usual arbitrage-pricing framework, we include in the market a
money market account \(B_t\), whose dynamics are now driven by the
stochastic short rate.

\paragraph{Definition.}
The money market account satisfies
\[
\mathrm{d}B_t=r_tB_t\,\mathrm{d}t,
\qquad
B_0=1.
\]

Solving this ordinary differential equation pathwise gives
\[
\boxed{
B_t=\exp\!\Big(\int_0^t r_u\,\mathrm{d}u\Big).
}
\]

\paragraph{Interpretation.}
The quantity \(B_t\) is the value at time \(t\) of one unit of cash
invested continuously at the short rate.

In the constant-rate case, this reduces to
\[
B_t=e^{rt}.
\]
So the money market account is the natural stochastic generalization of
the deterministic bank account used in Black--Scholes.

\paragraph{Discount factor.}
The associated discount factor between \(t\) and \(T\) is
\[
\frac{B_t}{B_T}
=
\exp\!\Big(-\int_t^T r_u\,\mathrm{d}u\Big).
\]

This is precisely the random discounting term that appears in the
risk-neutral pricing formula under stochastic rates.

\paragraph{Why this matters.}
In the Black--Scholes model with constant \(r\), discounting is easy
because
\[
\frac{B_t}{B_T}=e^{-r(T-t)}
\]
is deterministic.

Under stochastic rates, the discount factor itself is random, and this
changes both:
\begin{itemize}[leftmargin=1.4em]
    \item the pricing formula,
    \item and the structure of hedging portfolios.
\end{itemize}

\paragraph{Connection with zero-coupon bonds.}
The money market account is only one possible numéraire.
Another fundamental quantity in stochastic-rate models is the
zero-coupon bond price
\[
P(t,T),
\]
which gives the time-\(t\) value of one unit of currency paid at time
\(T\).

These bonds will play a central role later, because they are the natural
instruments used to hedge rate exposure and to perform changes of
numéraire.

\subsubsection{Examples of short-rate models}

The abstract dynamics
\[
\mathrm{d}r_t=b(t,r_t)\,\mathrm{d}t+\sigma_r(t,r_t)\,\mathrm{d}W_t
\]
become concrete once we choose a specific short-rate model.

We briefly recall three classical examples.

\paragraph{Vasicek model.}
In the Vasicek model, the short rate follows
\[
\boxed{
\mathrm{d}r_t=a(b-r_t)\,\mathrm{d}t+\sigma_r\,\mathrm{d}W_t.
}
\]

Here:
\begin{itemize}[leftmargin=1.4em]
    \item \(a>0\) is the mean-reversion speed,
    \item \(b\) is the long-run mean level,
    \item \(\sigma_r>0\) is the rate volatility.
\end{itemize}

\paragraph{Interpretation.}
This model is Gaussian and mean reverting.
It is analytically convenient, but it allows the short rate to become
negative.

\paragraph{Why it is useful pedagogically.}
The Vasicek model is often the simplest short-rate model for explicit
calculations.
It is therefore a good benchmark for theoretical discussions.

\paragraph{Cox--Ingersoll--Ross (CIR) model.}
In the CIR model, the short rate satisfies
\[
\boxed{
\mathrm{d}r_t=a(b-r_t)\,\mathrm{d}t+\sigma_r\sqrt{r_t}\,\mathrm{d}W_t.
}
\]

The square-root diffusion structure makes the model qualitatively
different from Vasicek.

\paragraph{Interpretation.}
Like Vasicek, the CIR model is mean reverting.
But because the volatility is proportional to \(\sqrt{r_t}\), the short
rate tends to remain non-negative under suitable parameter conditions.

This makes the model more realistic from an interest-rate perspective.

\paragraph{Hull--White model.}
A very common model in practice is the Hull--White model, which may be
written as
\[
\boxed{
\mathrm{d}r_t=\big(\theta(t)-a r_t\big)\,\mathrm{d}t+\sigma_r\,\mathrm{d}W_t.
}
\]

\paragraph{Interpretation.}
This is a time-inhomogeneous Gaussian short-rate model.
Its main practical advantage is that the function \(\theta(t)\) can be
chosen so as to fit the initial term structure exactly.

This makes Hull--White particularly attractive for practical calibration
and interest-rate derivative pricing.

\paragraph{What these models have in common.}
All these models aim to capture the same basic economic idea:
the short rate is random, but tends to revert toward some level or term
structure.

They differ mainly in:
\begin{itemize}[leftmargin=1.4em]
    \item whether rates can become negative,
    \item how volatility depends on the rate level,
    \item and how easily the model can be calibrated to the observed
    yield curve.
\end{itemize}

\paragraph{What matters for us in this section.}
For the pricing and hedging discussion that follows, the exact choice of
short-rate model is less important than the general consequences of
stochastic rates:
\begin{itemize}[leftmargin=1.4em]
    \item the option price depends on the current short rate,
    \item discounting becomes random,
    \item and rate-sensitive instruments will generally be needed for
    hedging.
\end{itemize}

So, in what follows, the reader may keep in mind either a generic
short-rate diffusion or a simple benchmark such as Vasicek or
Hull--White.

\subsection{Pricing a European option}

We now consider the pricing of a European option when the stock and the
short rate are both stochastic.

For definiteness, we take a payoff of the form
\[
H=g(S_T),
\]
for example
\[
g(S_T)=(S_T-K)^+.
\]

Even though the payoff depends only on the stock at maturity, the option
price will in general depend on the current short rate as well, because:
\begin{itemize}[leftmargin=1.4em]
    \item discounting is random,
    \item the stock drift under the risk-neutral measure is \(r_t\),
    \item and stock and rate shocks may be correlated.
\end{itemize}

Therefore, the pricing function is naturally written as
\[
u(t,s,r).
\]

\subsubsection{Risk-neutral pricing formula}

We begin with the fundamental risk-neutral valuation formula.

\paragraph{General pricing formula.}
Let \(H=g(S_T)\) be the payoff of a European option maturing at time
\(T\).
Under the risk-neutral measure \(\Qbb\), its price at time \(t\) is
\[
\boxed{
u(t,s,r)
=
\E^\Qbb\!\left[
\exp\!\Big(-\int_t^T r_u\,\mathrm{d}u\Big)\,g(S_T)
\;\middle|\;
S_t=s,\ r_t=r
\right].
}
\]

Equivalently, along the stochastic path,
\[
u(t,S_t,r_t)
=
\E^\Qbb\!\left[
\exp\!\Big(-\int_t^T r_u\,\mathrm{d}u\Big)\,g(S_T)
\;\middle|\;
\mathcal F_t
\right].
\]

\paragraph{Comparison with the constant-rate case.}
Under Black--Scholes with constant \(r\), the corresponding formula is
\[
u(t,s)
=
\E^\Qbb\!\left[
e^{-r(T-t)}g(S_T)
\;\middle|\;
S_t=s
\right].
\]

The only visible change is that the deterministic factor
\[
e^{-r(T-t)}
\]
has been replaced by the random discount factor
\[
\exp\!\Big(-\int_t^T r_u\,\mathrm{d}u\Big).
\]

But this apparently simple modification has major consequences:
\begin{itemize}[leftmargin=1.4em]
    \item the price now depends on the current short rate,
    \item the state space is enlarged,
    \item and the associated PDE becomes two-dimensional.
\end{itemize}

\paragraph{Why this formula is natural.}
The money market account satisfies
\[
B_t=\exp\!\Big(\int_0^t r_u\,\mathrm{d}u\Big),
\]
so the discounted value of any self-financing replicable claim must be a
martingale under \(\Qbb\).
This implies that the current option value is the conditional
expectation of the discounted payoff.

Thus the core principle of no-arbitrage pricing remains unchanged:
\begin{quote}
price equals discounted expectation under the risk-neutral measure.
\end{quote}

\paragraph{Markovian structure.}
Since we are working in a Markovian model for the pair
\[
(S_t,r_t),
\]
the option price depends only on the current time and the current values
of the two state variables:
\[
u=u(t,s,r).
\]

This is what allows us to derive a PDE in the extended state space.

\paragraph{A useful special case.}
For a European call with strike \(K\), the formula becomes
\[
\boxed{
u(t,s,r)
=
\E^\Qbb\!\left[
\exp\!\Big(-\int_t^T r_u\,\mathrm{d}u\Big)(S_T-K)^+
\;\middle|\;
S_t=s,\ r_t=r
\right].
}
\]

Even though the payoff depends only on \(S_T\), the price still depends
on \(r\) because the short rate affects both discounting and stock
dynamics.

\subsubsection{Pricing under a bond numeraire}

The risk-neutral formula above is correct, but it is not always the most
convenient one.
Since the payoff is paid at a fixed date \(T\), a more natural numéraire
is often the zero-coupon bond maturing at that same date.

\paragraph{Zero-coupon bond.}
Recall that
\[
P(t,T)
\]
denotes the time-\(t\) price of one unit of currency paid at time \(T\).

This asset is a natural numéraire for any claim with payment date \(T\).

\paragraph{Bond-numéraire pricing formula.}
Let \(\Qbb^T\) denote the \(T\)-forward measure associated with the bond
\(P(t,T)\).
Then the European option price can be written as
\[
\boxed{
u(t,s,r)
=
P(t,T)\,
\E^{\Qbb^T}\!\left[
g(S_T)
\;\middle|\;
S_t=s,\ r_t=r
\right].
}
\]

\paragraph{Why this is useful.}
This formula has a major advantage:
the random discount factor no longer appears explicitly inside the
expectation.
All discounting has been absorbed into the bond numéraire \(P(t,T)\).

So the pricing problem becomes:
\begin{itemize}[leftmargin=1.4em]
    \item multiply by the current bond price \(P(t,T)\),
    \item then take an expectation under the appropriate forward
    measure.
\end{itemize}

\paragraph{Comparison with the previous formula.}
Under the money market numéraire, we had
\[
u(t,s,r)
=
\E^\Qbb\!\left[
\exp\!\Big(-\int_t^T r_u\,\mathrm{d}u\Big)\,g(S_T)
\;\middle|\;
S_t=s,\ r_t=r
\right].
\]

Under the bond numéraire, this becomes
\[
u(t,s,r)
=
P(t,T)\,
\E^{\Qbb^T}\!\left[
g(S_T)
\;\middle|\;
S_t=s,\ r_t=r
\right].
\]

The two formulas are equivalent, but the second one is often more
convenient conceptually and computationally.

\paragraph{Interpretation.}
The \(T\)-forward measure is the probability measure under which prices
expressed in units of the bond \(P(t,T)\) are martingales.
This is the fixed-income analogue of the change-of-numéraire ideas
already seen earlier.

\paragraph{Example: European call.}
For a call payoff
\[
g(S_T)=(S_T-K)^+,
\]
the bond-numéraire formula reads
\[
\boxed{
u(t,s,r)
=
P(t,T)\,
\E^{\Qbb^T}\!\left[
(S_T-K)^+
\;\middle|\;
S_t=s,\ r_t=r
\right].
}
\]

So the option price may be viewed as:
\begin{itemize}[leftmargin=1.4em]
    \item the bond price \(P(t,T)\),
    \item times the forward-measure expectation of the payoff.
\end{itemize}

\paragraph{Main lesson.}
With stochastic rates, pricing can be written either:
\begin{itemize}[leftmargin=1.4em]
    \item under the risk-neutral measure with random discounting,
    \item or under a bond numéraire, with discounting absorbed into the
    bond price.
\end{itemize}

The second viewpoint is often cleaner and more natural for claims paid
at a fixed maturity.

\subsubsection{PDE in the extended state space}

We now derive the PDE satisfied by the pricing function
\[
u(t,s,r).
\]

\paragraph{Model dynamics.}
Assume that under the risk-neutral measure,
\[
\mathrm{d}S_t
=
r_tS_t\,\mathrm{d}t
+
\sigma_S(t,S_t,r_t)\,S_t\,\mathrm{d}W_t^{(1)},
\]
\[
\mathrm{d}r_t
=
b(t,r_t)\,\mathrm{d}t
+
\sigma_r(t,r_t)\,\mathrm{d}W_t^{(2)},
\]
with
\[
\mathrm{d}\langle W^{(1)},W^{(2)}\rangle_t=\rho\,\mathrm{d}t.
\]

Let \(u(t,s,r)\) be the price of a European claim with payoff
\[
u(T,s,r)=g(s).
\]

\paragraph{It\^o's formula.}
Applying It\^o's formula to \(u(t,S_t,r_t)\), we obtain
\begin{align*}
\mathrm{d}u(t,S_t,r_t)
&=
\partial_t u\,\mathrm{d}t
+
\partial_s u\,\mathrm{d}S_t
+
\partial_r u\,\mathrm{d}r_t
\\
&\quad
+
\frac12 \partial_{ss}u\,\mathrm{d}\langle S\rangle_t
+
\frac12 \partial_{rr}u\,\mathrm{d}\langle r\rangle_t
+
\partial_{sr}u\,\mathrm{d}\langle S,r\rangle_t.
\end{align*}

Now:
\[
\mathrm{d}\langle S\rangle_t
=
\sigma_S(t,S_t,r_t)^2 S_t^2\,\mathrm{d}t,
\]
\[
\mathrm{d}\langle r\rangle_t
=
\sigma_r(t,r_t)^2\,\mathrm{d}t,
\]
and
\[
\mathrm{d}\langle S,r\rangle_t
=
\rho\,\sigma_S(t,S_t,r_t)S_t\,\sigma_r(t,r_t)\,\mathrm{d}t.
\]

Substituting the dynamics of \(S_t\) and \(r_t\), we obtain
\begin{align*}
\mathrm{d}u(t,S_t,r_t)
&=
\Big(
\partial_t u
+
r_tS_t\,\partial_s u
+
b(t,r_t)\,\partial_r u
\\
&\qquad
+
\frac12 \sigma_S(t,S_t,r_t)^2 S_t^2\,\partial_{ss}u
+
\rho\,\sigma_S(t,S_t,r_t)S_t\,\sigma_r(t,r_t)\,\partial_{sr}u
\\
&\qquad
+
\frac12 \sigma_r(t,r_t)^2\,\partial_{rr}u
\Big)\mathrm{d}t
+\text{martingale terms.}
\end{align*}

\paragraph{Martingale condition for the discounted price.}
Under the risk-neutral measure, the discounted option price
\[
\frac{u(t,S_t,r_t)}{B_t}
\]
must be a martingale.
Therefore the drift of the discounted process must vanish.

This yields the PDE
\[
\boxed{
\partial_t u
+
rs\,\partial_s u
+
b(t,r)\,\partial_r u
+
\frac12 \sigma_S(t,s,r)^2 s^2\,\partial_{ss}u
+
\rho\,\sigma_S(t,s,r)s\,\sigma_r(t,r)\,\partial_{sr}u
+
\frac12 \sigma_r(t,r)^2\,\partial_{rr}u
-r\,u
=
0.
}
\]

The terminal condition is
\[
\boxed{
u(T,s,r)=g(s).
}
\]

\paragraph{Comparison with Black--Scholes.}
This is the natural extension of the Black--Scholes PDE to the case of
stochastic rates.

Compared with the constant-rate case, we now have:
\begin{itemize}[leftmargin=1.4em]
    \item an additional state variable \(r\),
    \item an additional first derivative \(\partial_r u\),
    \item an additional second derivative \(\partial_{rr}u\),
    \item and a mixed derivative \(\partial_{sr}u\) when stock and rate
    shocks are correlated.
\end{itemize}

\paragraph{Interpretation of the PDE.}
The terms in the PDE have a clear meaning:
\begin{itemize}[leftmargin=1.4em]
    \item \(rs\,\partial_s u\) comes from the stock drift under
    \(\Qbb\),
    \item \(b(t,r)\partial_r u\) comes from the short-rate drift,
    \item the second-order terms capture the local variances and
    covariance of the two factors,
    \item and the term \(-ru\) reflects discounting at the short rate.
\end{itemize}

\paragraph{A simple special case.}
If the stock volatility is constant and does not depend on the short
rate, so that
\[
\sigma_S(t,s,r)\equiv \sigma,
\]
then the PDE becomes
\[
\boxed{
\partial_t u
+
rs\,\partial_s u
+
b(t,r)\,\partial_r u
+
\frac12 \sigma^2 s^2\,\partial_{ss}u
+
\rho\,\sigma s\,\sigma_r(t,r)\,\partial_{sr}u
+
\frac12 \sigma_r(t,r)^2\,\partial_{rr}u
-r\,u
=
0.
}
\]

This is often the simplest form to keep in mind for a first treatment of
stochastic-rate option pricing.

\paragraph{Main message.}
With stochastic rates, the pricing PDE remains a no-arbitrage valuation
equation, but the problem now lives in the extended state space
\[
(t,s,r).
\]
So the conceptual structure is the same as in Black--Scholes, while the
dimension and the hedging implications become richer.

\subsection{Hedging with stochastic rates}

We now turn to the hedging problem in the presence of stochastic
interest rates.

The main new feature compared with the Black--Scholes benchmark is that
the option price depends not only on the stock price, but also on the
current short rate:
\[
u=u(t,s,r).
\]
Therefore, the option is exposed to both:
\begin{itemize}[leftmargin=1.4em]
    \item stock risk,
    \item and interest-rate risk.
\end{itemize}

This immediately suggests that a hedge based only on the stock will in
general not be sufficient.

\subsubsection{Why stock delta alone is not enough}

In the constant-rate Black--Scholes model, the option price is a
function
\[
u(t,s),
\]
and the hedge is naturally built from the stock and the money market
account.
The key quantity is the stock delta
\[
\Delta_t=\partial_s u(t,S_t).
\]

When rates are stochastic, however, the option price becomes
\[
u(t,s,r),
\]
so the current short rate \(r_t\) is now a state variable of the pricing
problem.

\paragraph{A first indication from the pricing function.}
Since the option price depends on \(r\), its total local variation
contains a term in \(\partial_r u\).
Applying It\^o's formula to \(u(t,S_t,r_t)\), one obtains diffusion
terms coming from both Brownian shocks:
\begin{align*}
\mathrm{d}u(t,S_t,r_t)
&=
\cdots
+
\partial_s u(t,S_t,r_t)\,\sigma_S(t,S_t,r_t)S_t\,\mathrm{d}W_t^{(1)}
\\
&\quad
+
\partial_r u(t,S_t,r_t)\,\sigma_r(t,r_t)\,\mathrm{d}W_t^{(2)}.
\end{align*}

So the option is locally exposed to two sources of randomness.

\paragraph{Why the stock cannot hedge rate risk by itself.}
If we hedge only with the stock and cash, then the stock position can
only generate exposure to the Brownian shocks carried by the stock
process itself.
It does not automatically span an independent short-rate shock.

Therefore, even if we choose the stock holding so as to match the stock
sensitivity
\[
\partial_s u,
\]
we will generally still be left with residual exposure to the short-rate
factor through \(\partial_r u\).

\paragraph{A simple one-factor intuition.}
Suppose for a moment that the stock and the short rate are driven by two
independent Brownian motions.
Then a hedge based only on:
\begin{itemize}[leftmargin=1.4em]
    \item the stock,
    \item and the money market account,
\end{itemize}
can neutralize the stock risk, but not the independent rate risk.

So, unless the rate factor is degenerate or fully spanned by the stock,
stock delta alone cannot produce exact replication.

\paragraph{What this means financially.}
An option on the stock may be exposed to rate movements even if its
payoff depends only on \(S_T\).
This happens because:
\begin{itemize}[leftmargin=1.4em]
    \item discounting depends on rates,
    \item the risk-neutral drift of the stock depends on rates,
    \item and stock and rate shocks may be correlated.
\end{itemize}

Thus, under stochastic rates, a hedge based only on the stock is in
general incomplete.

\paragraph{Main lesson.}
The classical stock delta remains an important sensitivity, but it is no
longer the whole story.
Once rates are stochastic, the option carries a second source of risk,
and hedging requires rate-sensitive instruments.

\subsubsection{Hedging with stock and bonds}

A natural way to hedge rate exposure is to include in the hedging
portfolio a traded fixed-income instrument, for example a zero-coupon
bond.

\paragraph{Choice of hedging instruments.}
We consider a hedging portfolio built from:
\begin{itemize}[leftmargin=1.4em]
    \item the stock \(S_t\),
    \item a zero-coupon bond \(P(t,T^\star)\) with some maturity
    \(T^\star\),
    \item and the money market account \(B_t\).
\end{itemize}

We write the portfolio value as
\[
\Pi_t=\phi_t S_t+\psi_t P(t,T^\star)+\beta_t B_t.
\]

Here:
\begin{itemize}[leftmargin=1.4em]
    \item \(\phi_t\) is the stock position,
    \item \(\psi_t\) is the bond position,
    \item \(\beta_t\) is the cash position.
\end{itemize}

\paragraph{Why a bond is the right additional instrument.}
A zero-coupon bond is directly sensitive to the short rate.
Its price moves when the interest-rate environment changes.
So it is the natural instrument for hedging the rate exposure of the
option.

This is the stochastic-rate analogue of using the stock to hedge stock
risk.

\paragraph{Dynamics of a zero-coupon bond under the risk-neutral measure.}
Let \(P(t,T)\) denote the price at time \(t\) of a zero-coupon bond
maturing at time \(T\), and let \(B_t\) be the money market account,
defined by
\[
\mathrm{d}B_t=r_t B_t\,\mathrm{d}t.
\]
Under the risk-neutral measure associated with the money market account,
the discounted bond price
\[
\frac{P(t,T)}{B_t}
\]
must be a martingale. Therefore, the bond price has drift
\(r_tP(t,T)\). In a diffusion setting, its dynamics can be written as
\[
\mathrm{d}P(t,T)
=
r_t P(t,T)\,\mathrm{d}t
+
P(t,T)\,\Sigma_P(t,T)\cdot \mathrm{d}W_t^{\Qbb},
\]
or equivalently
\[
\frac{\mathrm{d}P(t,T)}{P(t,T)}
=
r_t\,\mathrm{d}t
+
\Sigma_P(t,T)\cdot \mathrm{d}W_t^{\Qbb},
\]
where \(\Sigma_P(t,T)\) denotes the volatility vector of the bond price.
Thus, under the risk-neutral measure, a zero-coupon bond behaves like
any traded asset: its drift is the short rate, while its volatility
depends on the maturity \(T\). In particular, the process
\[
\frac{P(t,T)}{B_t}
\]
is a martingale.

In a short-rate model, bond prices are typically obtained as functions
of the current short rate, namely
\[
P(t,T)=F(t,r_t).
\]
Applying It\^o's formula then shows that the bond volatility is directly
induced by the volatility of the short rate. More precisely, if
\[
\mathrm{d}r_t=b(t,r_t)\,\mathrm{d}t+\sigma_r(t,r_t)\,\mathrm{d}W_t,
\]
then the diffusion part of \(P(t,T)\) is given by
\[
\partial_r F(t,r_t)\,\sigma_r(t,r_t)\,\mathrm{d}W_t.
\]
So the volatility of the zero-coupon bond is directly linked to two
quantities:
\begin{itemize}[leftmargin=1.4em]
    \item the volatility of the short rate itself,
    \item and the sensitivity of the bond price to the short rate.
\end{itemize}
In this sense, short-rate models generate bond-price dynamics
endogenously: once the process \(r_t\) is specified, the corresponding
zero-coupon bond dynamics follow. Conversely, one may also choose to
model bond prices or forward rates directly, as in HJM-type frameworks,
rather than starting from the short rate.

\paragraph{Local replication idea.}
Suppose that the stock and the bond have diffusive dynamics of the form
\[
\mathrm{d}S_t=\cdots+\Gamma_S^{(1)}\,\mathrm{d}W_t^{(1)}+\Gamma_S^{(2)}\,\mathrm{d}W_t^{(2)},
\]
\[
\mathrm{d}P(t,T^\star)=\cdots+\Gamma_P^{(1)}\,\mathrm{d}W_t^{(1)}+\Gamma_P^{(2)}\,\mathrm{d}W_t^{(2)}.
\]

The option price has diffusion terms
\[
\mathrm{d}u(t,S_t,r_t)
=
\cdots+\Gamma_u^{(1)}\,\mathrm{d}W_t^{(1)}+\Gamma_u^{(2)}\,\mathrm{d}W_t^{(2)}.
\]

If the matrix
\[
\begin{pmatrix}
\Gamma_S^{(1)} & \Gamma_S^{(2)}\\
\Gamma_P^{(1)} & \Gamma_P^{(2)}
\end{pmatrix}
\]
has full rank, then we can choose \(\phi_t\) and \(\psi_t\) so as to
match both Brownian exposures of the option.

This is the local replication condition.

\paragraph{Interpretation.}
The stock hedges the stock-like part of the option's exposure.
The bond hedges the rate-like part.
The cash account then adjusts the deterministic part of the portfolio.

So, in the simplest two-factor setting, the natural analogue of the
Black--Scholes delta hedge becomes:
\begin{quote}
stock \(+\) bond \(+\) cash.
\end{quote}

\paragraph{The role of the bond maturity.}
The maturity \(T^\star\) of the hedging bond matters.
Different bonds have different sensitivities to the short rate and to the
yield curve.

In a one-factor short-rate model, a single bond of suitable maturity may
already be enough to span the rate risk.
In richer term-structure models, one may need several bonds or swaps to
hedge all relevant rate factors.

\paragraph{A useful special case.}
Suppose the market has only two Brownian sources of risk:
\begin{itemize}[leftmargin=1.4em]
    \item one stock-related factor,
    \item one rate-related factor.
\end{itemize}

If the stock and the chosen bond span these two directions, then a
European option can be replicated locally using:
\[
\phi_t \text{ units of stock}
\quad + \quad
\psi_t \text{ units of bond}
\quad + \quad
\beta_t \text{ in cash.}
\]

So stochastic rates do not automatically imply incompleteness.
Completeness depends on whether the traded assets span all the sources of
risk.

\paragraph{Practical viewpoint.}
Even when exact replication is not available globally, the same idea is
still useful:
\begin{itemize}[leftmargin=1.4em]
    \item use the stock to hedge stock exposure,
    \item use bonds or swaps to hedge rate exposure,
    \item and interpret the residual error as model or spanning error.
\end{itemize}

Thus the bond-augmented hedge is the natural generalization of the
classical Black--Scholes hedge.

\subsubsection{Sensitivity to the short rate}

Since the price depends on the current short rate, a natural quantity to
consider is the rate sensitivity
\[
\partial_r u(t,s,r).
\]

\paragraph{Definition.}
The quantity
\[
\partial_r u(t,s,r)
\]
measures how the option price changes when the current short rate is
perturbed, all else being held fixed.

It is the analogue, in the rate direction, of the stock delta
\[
\partial_s u(t,s,r).
\]

\paragraph{Why this sensitivity appears.}
The short rate affects the option price through several channels:
\begin{itemize}[leftmargin=1.4em]
    \item it changes the discount factor,
    \item it changes the risk-neutral drift of the stock,
    \item and it may be correlated with the stock shocks.
\end{itemize}

So even if the payoff depends only on \(S_T\), the current rate matters
for the current value of the option.

\paragraph{Relation with fixed-income hedging.}
In practice, the quantity \(\partial_r u\) is not usually hedged by
trading directly in the short rate, since the short rate is not itself a
traded asset.
Instead, one hedges this rate sensitivity using traded fixed-income
instruments such as:
\begin{itemize}[leftmargin=1.4em]
    \item zero-coupon bonds,
    \item bond futures,
    \item swaps,
    \item or swaptions in richer settings.
\end{itemize}

Thus the short-rate sensitivity should be understood as the source of
the hedge, not as the traded hedge instrument itself.

\paragraph{A simple interpretation in PDE terms.}
The presence of the derivative \(\partial_r u\) and the second-order
term \(\partial_{rr}u\) in the pricing PDE already tells us that the
option price reacts to movements in the short rate.
So the PDE itself makes visible the need for an interest-rate hedge.

\paragraph{Comparison with the constant-rate case.}
In Black--Scholes with constant \(r\), one may still differentiate with
respect to the parameter \(r\), but the rate is not a state variable.
So such a sensitivity is only a parameter sensitivity.

With stochastic rates, by contrast, \(r_t\) is a genuine random state
variable.
Therefore \(\partial_r u\) is a true dynamic sensitivity of the option
price.

\paragraph{Main message.}
Under stochastic rates, the option carries not only stock delta risk but
also short-rate risk.
This is why the natural hedge is no longer built solely from the stock:
one also needs fixed-income instruments to manage the sensitivity of the
price to the current rate environment.

\subsection{Replication versus incompleteness}

We now address one of the central conceptual questions raised by
stochastic interest-rate models:

\begin{quote}
Can the option payoff still be replicated exactly, as in the standard
Black--Scholes setting, or does the market become incomplete?
\end{quote}

The answer depends on the number of underlying sources of risk and on
the set of traded assets available for hedging.

This is an important point, because it determines whether the pricing
problem can be understood as an exact replication problem or only as a
problem of approximate hedging.

\subsubsection{When exact replication is possible}

The guiding principle is the following:
\begin{quote}
exact replication is possible when the traded assets span all sources of
risk driving the option price.
\end{quote}

\paragraph{A reminder from Black--Scholes.}
In the standard Black--Scholes model, there is only one Brownian source
of uncertainty.
The risky asset \(S_t\) carries this source of risk, while the money
market account is deterministic once the short rate is fixed.

Therefore, one risky asset together with the money market account is
enough to replicate a European option.
This is why the model is complete.

\paragraph{What changes with stochastic rates.}
When the short rate is stochastic, the state variable becomes
\[
(t,S_t,r_t),
\]
and there may now be two sources of randomness:
\begin{itemize}[leftmargin=1.4em]
    \item one affecting the stock,
    \item one affecting the short rate.
\end{itemize}

If the option price depends on both \(S_t\) and \(r_t\), then exact
replication requires enough traded risky assets to hedge both sources of
risk.

\paragraph{A simple two-factor situation.}
Suppose that under the risk-neutral measure the economy is driven by two
Brownian motions and that we can trade:
\begin{itemize}[leftmargin=1.4em]
    \item the stock \(S_t\),
    \item a zero-coupon bond \(P(t,T^\star)\),
    \item and the money market account \(B_t\).
\end{itemize}

If the stock and the bond have linearly independent exposures to the two
Brownian sources, then they span the whole risk space.

In that case, a self-financing portfolio of the form
\[
\Pi_t=\phi_t S_t+\psi_t P(t,T^\star)+\beta_t B_t
\]
can be chosen so as to match both diffusion terms of the option price.

\paragraph{Local replication condition.}
This is most clearly seen at the level of the diffusion coefficients.
If the stock and the bond satisfy
\[
\mathrm{d}S_t=\cdots+\Gamma_S^{(1)}\,\mathrm{d}W_t^{(1)}+\Gamma_S^{(2)}\,\mathrm{d}W_t^{(2)},
\]
\[
\mathrm{d}P(t,T^\star)=\cdots+\Gamma_P^{(1)}\,\mathrm{d}W_t^{(1)}+\Gamma_P^{(2)}\,\mathrm{d}W_t^{(2)},
\]
and the option price satisfies
\[
\mathrm{d}u(t,S_t,r_t)=\cdots+\Gamma_u^{(1)}\,\mathrm{d}W_t^{(1)}+\Gamma_u^{(2)}\,\mathrm{d}W_t^{(2)},
\]
then exact local replication is possible if the linear system
\[
\phi_t \Gamma_S^{(1)}+\psi_t \Gamma_P^{(1)}=\Gamma_u^{(1)},
\]
\[
\phi_t \Gamma_S^{(2)}+\psi_t \Gamma_P^{(2)}=\Gamma_u^{(2)}
\]
admits a unique solution.

This is the case precisely when the matrix
\[
\begin{pmatrix}
\Gamma_S^{(1)} & \Gamma_P^{(1)}\\
\Gamma_S^{(2)} & \Gamma_P^{(2)}
\end{pmatrix}
\]
is invertible.

\paragraph{Interpretation.}
So, in a two-factor model:
\begin{itemize}[leftmargin=1.4em]
    \item if we have two risky traded assets with independent exposures,
    \item then exact replication may still be possible.
\end{itemize}

Thus, stochastic rates do \emph{not} automatically imply market
incompleteness.
What matters is whether the available traded assets span all the
relevant risk factors.

\paragraph{Examples where replication may remain possible.}
In a one-factor short-rate model, if both the stock and all bond prices
are driven by the same single Brownian motion, then the market may still
be complete.

More generally, in a two-factor model, exact replication may be
available if we can trade the stock and a sufficiently sensitive
interest-rate instrument.

So the correct message is:
\[
\boxed{
\text{stochastic rates enlarge the problem, but do not automatically
destroy replication.}
}
\]

\subsubsection{When the market is incomplete}

Although exact replication may be possible in some models, it is not
automatic.
In many situations, stochastic-rate models lead naturally to market
incompleteness.

\paragraph{Basic source of incompleteness.}
The market becomes incomplete whenever there are more relevant sources of
risk than traded assets available to hedge them.

For example, suppose that:
\begin{itemize}[leftmargin=1.4em]
    \item the stock and the short rate are driven by two independent
    Brownian factors,
    \item but the only traded risky asset is the stock itself.
\end{itemize}

Then a portfolio built from:
\begin{itemize}[leftmargin=1.4em]
    \item the stock,
    \item and the money market account,
\end{itemize}
cannot hedge the rate factor independently.
So exact replication fails.

\paragraph{Why the money market account does not solve the problem.}
This is an important point.
The money market account evolves as
\[
\mathrm{d}B_t=r_tB_t\,\mathrm{d}t,
\]
so its local variation contains no independent Brownian term.
It is therefore not a risky asset in the sense needed to span a new
diffusion direction.

Thus adding the bank account does not add a new hedging factor.
To hedge a genuinely stochastic rate factor, one needs a traded
rate-sensitive risky instrument, such as a bond or a swap.

\paragraph{More general sources of incompleteness.}
Even if one can trade a stock and some bonds, the market may still be
incomplete if:
\begin{itemize}[leftmargin=1.4em]
    \item the term-structure model is driven by more factors than the
    available traded instruments span,
    \item volatility is itself stochastic,
    \item or there are additional sources of basis risk or model risk.
\end{itemize}

So incompleteness is not a special pathology of one particular model.
It is the natural outcome whenever the risk structure is richer than the
traded asset set.

\paragraph{Consequence for pricing.}
When exact replication is not possible, no-arbitrage alone does not
necessarily pin down a unique price via replication.
Instead, pricing depends on the chosen martingale measure or, more
practically, on the chosen pricing model.

Thus the link between
\[
\text{price} \quad \text{and} \quad \text{exact replication}
\]
becomes weaker than in the Black--Scholes benchmark.

\paragraph{Consequence for hedging.}
In an incomplete market, the best one can generally do is:
\begin{itemize}[leftmargin=1.4em]
    \item hedge the main risk factors,
    \item minimize some criterion for the residual error,
    \item and interpret the remaining exposure as unhedgeable risk.
\end{itemize}

So the hedging problem becomes one of risk management rather than exact
replication.

\paragraph{Main message.}
Under stochastic rates, incompleteness appears whenever the option is
exposed to risk factors that are not fully spanned by the traded asset
set.
In that case, the option can still be priced and hedged, but not
replicated perfectly.

\subsubsection{Practical hedging interpretation}

From a practical perspective, the distinction between complete and
incomplete markets changes the way one interprets hedging.

\paragraph{In a complete model.}
If the market is complete, then the hedging strategy has a strong
interpretation:
\begin{itemize}[leftmargin=1.4em]
    \item it replicates the payoff exactly,
    \item it determines the option price uniquely by no-arbitrage,
    \item and any deviation from the model hedge is a genuine
    mis-hedge.
\end{itemize}

This is the idealized Black--Scholes picture.

\paragraph{In an incomplete model.}
If the market is incomplete, then the hedge no longer has the status of
an exact replication strategy.
Instead, it should be viewed as an \emph{approximate} hedge designed to
control the main sources of risk.

In practice, one typically:
\begin{itemize}[leftmargin=1.4em]
    \item hedges the stock sensitivity with the stock,
    \item hedges the rate sensitivity with bonds, swaps, or related
    instruments,
    \item and accepts a residual hedging error that cannot be eliminated
    exactly.
\end{itemize}

\paragraph{What the sensitivities mean in practice.}
The stock delta
\[
\partial_s u(t,s,r)
\]
still measures the local exposure to the stock.
The rate sensitivity
\[
\partial_r u(t,s,r)
\]
measures the local exposure to the short-rate environment.

These sensitivities guide the construction of a hedge, even when exact
replication is impossible.

\paragraph{Practical objective.}
So, in practice, the aim is often not:
\[
\text{``replicate exactly.''}
\]
Rather, it is:
\[
\text{``hedge the dominant factors and reduce the residual risk as much
as possible.''}
\]

This is a much more realistic interpretation of hedging in multi-factor
markets.

\paragraph{Why the distinction still matters pedagogically.}
Even if exact replication is not always available in practice, the
complete-market benchmark remains extremely useful:
\begin{itemize}[leftmargin=1.4em]
    \item it clarifies what a perfect hedge would mean,
    \item it helps us understand the structure of the PDE,
    \item and it shows which additional instruments are needed to span
    the risk.
\end{itemize}

So the theory of replication remains central even when the real market
is only approximately hedgeable.

\paragraph{Main practical message.}
With stochastic rates, one should think of hedging as a layered problem:
\begin{itemize}[leftmargin=1.4em]
    \item stock hedge for stock risk,
    \item bond or swap hedge for rate risk,
    \item and residual exposure if the market does not span all factors.
\end{itemize}

Thus the presence of stochastic rates enriches both the theory and the
practice of hedging.

\subsection{Main lessons}

This section extends the standard option-pricing logic to the case of
stochastic interest rates.

The central ideas are the following.

\paragraph{First lesson: random discounting changes pricing.}
Once the short rate becomes stochastic, the discount factor is no longer
deterministic.
As a result, even a payoff depending only on \(S_T\) leads to a price of
the form
\[
u(t,s,r)
=
\E^\Qbb\!\left[
\exp\!\Big(-\int_t^T r_u\,\mathrm{d}u\Big)\,g(S_T)
\;\middle|\;
S_t=s,\ r_t=r
\right].
\]

So the option price depends naturally on both the stock price and the
current short rate.

\paragraph{Second lesson: the state space is enlarged.}
With stochastic rates, the natural state variable is
\[
(t,S_t,r_t),
\]
and the pricing PDE becomes two-dimensional in the spatial variables.
This leads to additional first- and second-order terms, including a
mixed derivative term when stock and rate shocks are correlated.

\paragraph{Third lesson: a bond numéraire is often the natural viewpoint.}
For a payoff paid at time \(T\), it is often convenient to use the
zero-coupon bond \(P(t,T)\) as numéraire.
This leads to the formula
\[
u(t,s,r)
=
P(t,T)\,
\E^{\Qbb^T}\!\left[
g(S_T)
\;\middle|\;
S_t=s,\ r_t=r
\right],
\]
which is often cleaner than working directly with random discounting
under the money market numéraire.

\paragraph{Fourth lesson: stock delta alone is no longer enough.}
Since the option price depends on the short rate, the hedge must account
for both:
\begin{itemize}[leftmargin=1.4em]
    \item stock exposure,
    \item and rate exposure.
\end{itemize}

Therefore, the natural hedge is not only a stock hedge, but generally a
combination of:
\begin{itemize}[leftmargin=1.4em]
    \item stock,
    \item fixed-income instruments such as bonds or swaps,
    \item and cash.
\end{itemize}

\paragraph{Fifth lesson: replication depends on spanning.}
Stochastic rates do not automatically imply incompleteness.
Exact replication is possible whenever the traded risky assets span all
relevant sources of randomness.

But if the market does not span all factors, then the option cannot be
replicated perfectly, and hedging becomes approximate rather than exact.

\paragraph{Sixth lesson: hedging becomes a risk-management problem.}
In incomplete markets, hedging should be interpreted as the management
of the dominant exposures, not as exact replication.
The residual error is part of the problem and must be understood as
unhedgeable risk.

\paragraph{Final summary.}
So the general philosophy remains the same as before:
\[
\text{model}
\longrightarrow
\text{risk-neutral pricing formula}
\longrightarrow
\text{PDE}
\longrightarrow
\text{hedging strategy}.
\]

But with stochastic rates, every step becomes richer:
\begin{itemize}[leftmargin=1.4em]
    \item discounting is random,
    \item the state space is larger,
    \item the hedge requires more instruments,
    \item and exact replication is no longer guaranteed.
\end{itemize}


\end{document}